%! TEX root = main.tex
% Stile del documento
\documentclass[12pt,openany]{amsbook}
\usepackage[T1]{fontenc}
\usepackage[utf8]{inputenc}
\usepackage[a4paper]{geometry}
\geometry{verbose,tmargin=1.8cm,bmargin=2cm,lmargin=2cm,rmargin=2cm,headheight=1.8cm,headsep=1cm,footskip=1cm}
% Localizzazione (traduce nomi automatici nella lingua selezionata)
\usepackage[italian]{babel}
% Altri pacchetti per la formattazione
\usepackage{lmodern}
\usepackage{amsthm}
\usepackage{amsmath}
\usepackage{mathrsfs}
\usepackage{mathtools}
\usepackage{cancel}
\usepackage{enumerate}
\usepackage{graphicx}
 % Impostazioni dei metadati del PDF
\usepackage[unicode=true, bookmarks=true, bookmarksnumbered=true, bookmarksopen=true, bookmarksopenlevel=1, breaklinks=false, pdfborder={0 0 0}, pdfborderstyle={}, backref=page, colorlinks=false]
{hyperref}
\hypersetup{
    pdftitle={Appunti di Probbilità e Statistica},
    pdfauthor={Lorenzo Prosseda},
    pdfsubject={Corso di Probabilità e Statistica del prof. Barchielli, Politecnico di Milano 2018\---2019},
    pdfkeywords={probabilità, statistica}
}
% Percorso della cartella contenente le immagini
\graphicspath{ {../img/} }
% Impostazioni di stile per gli ambienti teorema, definizione, esempio, ecc
\theoremstyle{plain}
\newtheorem{thm}{Teorema}[section] % resetta numerazione ad ogni capitolo
\renewcommand{\thethm}{\arabic{chapter}.\arabic{section}.\arabic{thm}} % numerazione ambienti teorema
\numberwithin{equation}{section} % numerazione equazioni
% Definizione ambienti personalizzati
\theoremstyle{definition}
\newtheorem{defn}[thm]{Definizione}
\newtheorem{prty}[thm]{Proprietà}
\newtheorem{exmp}[thm]{Esempio}
\newtheorem{obsv}[thm]{Osservazione}
\theoremstyle{remark}
\newtheorem{note}[thm]{Nota}
% Numerazione personalizzata sezioni e sottosezioni
\renewcommand\thesection{\thechapter.\arabic{section}}
\renewcommand\thesubsection{\thesection.\arabic{subsection}}
% Comando trattino "non spezzante" evita sillabazione sul trattino
\newcommand\nbdash{\nobreakdash-\hspace{0pt}}

% In questo ambiente risiede il contenuto del documento
\begin{document}
    % Pagina del titolo
    %! TEX root = main.tex

\begin{titlepage}
    \centering % centra tutto in questo ambiente
    \vspace*{\stretch{0.8}}
    \Huge\textbf{Appunti di Probabilità e Statistica}\\
    \vspace*{\stretch{0.6}}
    \Large\textit{Lorenzo Prosseda}\\
    \bigskip
    \large a.a. 2018\---2019 \\
    \vspace*{\stretch{5.0}}
    \begin{figure}[h] % inserisci logo ed estratto della licenza GFDL
        \includegraphics[width=2cm]{gfdl-logo.pdf}
        \bigskip \\
        \small\textit{Copyright \textcopyright\,2019 Lorenzo Prosseda. Permission is granted to copy, distribute and/or modify this document under the terms of the GNU Free Documentation License, Version 1.3 or any later version published by the Free Software Foundation; with no Invariant Sections, no Front-Cover Texts, and no Back-Cover Texts. A copy of the license is included in the file called ``LICENSE''.}
    \end{figure}
    \vspace*{\stretch{1.0}}
\end{titlepage}


    % Genera indice di capitoli e sezioni
    \tableofcontents

    % Includi file dei capitoli
    \includeonly{chap3}
    %! TEX root = main.tex
% Capitolo 1

\chapter{Teoria della Probabilità}
    \section{Spazio dei campioni}
        \begin{defn}
            Sia dato un esperimento \emph{aleatorio} (impossibile prevederne risultato), i cui risultati siano rappresentati da $\omega \in \Omega$; chiamiamo l'insieme $\Omega$:
            \begin{itemize}
                \item spazio campionario;
                \item spazio dei campioni;
                \item spazio degli eventi \underline{elementari};
                \item spazio degli esiti;
            \end{itemize}
            diciamo inoltre che $\Omega$ è relativo all'esperimento effettuato, e chiamiamo i suoi elementi $\omega$ \emph{eventi \underline{elementari}}.
        \end{defn}
        \section{Eventi}\label{sec:Eventi}
        \begin{defn}
            I sottoinsiemi dello spazio campionario e le loro combinazioni (in termini di eventi elementari) tramite operatori logici di unione ($\cup$), intersezione ($\cap$) e negazione ($^{\text{C}}$) sono chiamati \emph{eventi} (non sono più elementari).
        \end{defn}
        \begin{obsv}
            Gli eventi \emph{elementari} sono rappresentati da sottoinsiemi di $\Omega$ di cardinalità 1: un evento elementare $E$ è definito come $\{\omega_E\} \in \Omega \land |\{\omega_E\}| = 1$.
        \end{obsv}
        \begin{defn}
            Gli eventi $E_i$ possono essere rappresentati con sottoinsiemi dello spazio campionario $\Omega$, dunque essi formano una \emph{famiglia} o \emph{collezione} di sottoinsiemi di $\Omega$, che indichiamo con $\mathscr{F}$.

            Questo implica che un evento che faccia parte di questa famiglia di sottoinsiemi sia contenuto nello spazio campionario ($E \in \mathscr{F} \implies E \subset \Omega$).

            Diremo che si è \emph{verificato} un evento $E$ nel contesto di un esperimento aleatorio, se almeno uno degli esiti è contenuto in esso ($\omega \in E$).
        \end{defn}
        \begin{defn}
            Dati due eventi $E$ ed $F$, sono definite le operazioni di unione $E \cup F$ (gli eventi elementari in $E$ o in $F$) e intersezione $E \cap F$ (gli eventi elementari presenti sia in $E$ che in $F$).
            \begin{itemize}
                \item Gli eventi che non contengono alcun evento elementare sono chiamati \emph{eventi vuoti} e si indicano con \O. Se due eventi non hanno eventi elementari in comune ($E \cap F = \text{\O}$), essi si dicono \emph{disgiunti}.
                \item Per ogni evento $E \in \mathscr{F}$ è definito il complementare $E^{\text{C}}$ come l'insieme degli eventi elementari di $\mathscr{F}$ che non stanno in $E$.
                \item Se tutti gli eventi elementari di un evento $E$ sono anche in un evento  $F$, diremo che $E$ \emph{è contenuto} in  $F$ ($E \subset F$); se vale anche l'inverso $E \supset F$ allora diremo che i due eventi sono \emph{uguali} ($E \equiv F$).
            \end{itemize}
            Indichiamo l'unione di $n$ eventi con  $\bigcup_{k=1}^{n} E_k$ e l'intersezione di  $n$ eventi con $\bigcap_{k=1}^{n} E_k$.
        \end{defn}
        \begin{defn}
            Sia $\Omega$ uno spazio campionario e $\mathscr{F}$ una famiglia di sottoinsiemi di $\Omega$; diremo che $\mathscr{F}$ è un'\emph{algebra di sottoinsiemi} di $\Omega$ se soddisfa le seguenti proprietà:
            \begin{enumerate}
                \item $\Omega \in \mathscr{F}$;
                \item $E \in \mathscr{F} \implies E^{\text{C}} \coloneqq \Omega \backslash E \in \mathscr{F}$;
                \item $E,\,F \in \mathscr{F} \implies \{E \cup F\} \in \mathscr{F}$.
            \end{enumerate}
        \end{defn}
        \begin{defn}
            Sia $\Omega$ uno spazio campionario e $\mathscr{F}$ una famiglia di sottoinsiemi di $\Omega$; diremo che $\mathscr{F}$ è una \emph{$\sigma$\nbdash algebra di sottoinsiemi} di $\Omega$ se soddisfa le seguenti proprietà:
            \begin{enumerate}
                \item $\Omega \in \mathscr{F}$;
                \item $E \in \mathscr{F} \implies E^{\text{C}} \coloneqq \Omega \backslash E \in \mathscr{F}$;
                \item $E_1,\,E_2,\,\ldots \in \mathscr{F} \implies \bigcup_{k=1}^{+\infty}E_k \in \mathscr{F}$.
            \end{enumerate}
        \end{defn}
        \begin{defn}
            La coppia di spazio campionario e famiglia di suoi sottoinsiemi ($\sigma$\nbdash algebra) $(\Omega,\,\mathscr{F})$ è chiamata \emph{spazio probabilizzabile}.
        \end{defn}
    \section{Spazio di probabilità}
        \begin{defn}\label{defn:Probabilità_impostazione_assiomatica}
            Sia dato uno spazio probabilizzabile $(\Omega,\,\mathscr{F})$; chiamiamo \emph{probabilità} su $(\Omega,\,\mathscr{F})$ una funzione $P$ su $\mathscr{F}$ tale che:
            \begin{enumerate}
                \item $\forall E \in \mathscr{F}\,:\,P(E) \geq 0$;
                \item $P(\Omega) = 1$;
                \item $\forall h,k\,:\,h \neq k \implies \forall E \in \mathscr{F}\,:\,E_h \cap E_k = \text{\O}$ ($\sigma$\nbdash additività o \emph{additività completa}).
            \end{enumerate}
        La terna $(\Omega,\,\mathscr{F},\,P)$, con $\mathscr{F}$ $\sigma$\nbdash algebra, si chiama \emph{spazio di probabilità} (impostazione assiomatica).
        \end{defn}
        \begin{prty}\label{prty:Spazio_di_probabilità}
            Sia $(\Omega,\,\mathscr{F},\,P)$ uno spazio di probabilità; allora vale:
            \begin{enumerate}
                \item $P(\text{\O}) = 0$ (\emph{evento impossibile});
                \item $\forall h,k\,:\,h \neq k \implies \forall E \in \mathscr{F}\,:\,E_h \cap E_k = \text{\O} \implies P(\bigcup_{k=1}^{n} E_k) = \sum_{k=1}^{n} P(E_k)$ (\emph{additività finita}).
            \end{enumerate}
        \end{prty}
        \begin{proof}
            \hfill
            \begin{enumerate}
                \item $\forall k \in [1,\,n]\,:\,E_k \coloneqq \text{\O} \implies \bigcup_{k=1}^{+\infty} E_k = \text{\O}$, inoltre $E_1,\,\ldots,\,E_n$ è una successione di eventi disgiunti a coppie; per l'assioma $(3)$ della Definizione~\ref{defn:Probabilità_impostazione_assiomatica} vale inoltre: \[
                        P\left(\text{\O}\right) = P\left(\bigcup_{k=1}^{+\infty} E_k\right) = \sum_{k=1}^{+\infty} P\left(\text{\O}\right)
                .\] 
            \item Se $\forall k \in [n+1,\,\infty)\,:\,E_k = \text{\O}$, allora $E_1,\,E_2,\,\ldots$ è una successione di eventi disgiunti a coppie, e vale $\bigcup_{k=1}^{+\infty} E_k = \bigcup_{k=1}^{n} E_k$; per l'assioma $(3)$ della Definizione~\ref{defn:Probabilità_impostazione_assiomatica} vale inoltre: \[
                    P(\bigcup_{k=1}^{n} E_k) = (\bigcup_{k=1}^{+\infty} E_k) = \sum_{k=1}^{n} E_k + \underset{E_k = \text{\O}}{\underbrace{\cancel{\sum_{k=n+1}^{+\infty} E_k}}} = \sum_{k=1}^{n} E_k
            . \qedhere\]
            \end{enumerate}
        \end{proof}
    \section{Definizione assiomatica di Probabilità}
        \begin{defn}
            Se assegnamo a ogni evento in uno spazio di probabilità $E_k \in (\Omega,\,\mathscr{F},\,P)$ un valore $x = P(E_k)$ che chiamiamo probabilità dell'evento $E$, esso deve rispettare i seguenti assiomi, dedotti dalla Definizione~\ref{defn:Probabilità_impostazione_assiomatica}:
            \begin{enumerate}[\indent I)]
                \item $0 \leq P(E_k) \leq 1$;
                \item $P(\Omega) = 1$;
                \item  se gli eventi dello spazio di probabilità sono disgiunti a coppie otteniamo: \[
                \forall n \in [1,\,\infty)\,:\,P(\bigcup_{k=1}^{n} E_k) = \sum_{k=1}^{n} P(E_k)
                .\] 
            \end{enumerate}
        \end{defn}
    \section{Proprietà della funzione Probabilità}
    \begin{prty}\label{prty:Proprietà_funzione_probabilità}
            Sia $(\Omega,\,\mathscr{F},\,P)$ uno spazio di probabilità; su di esso possiamo osservare le seguenti proprietà:
            \begin{enumerate}
                \item $E \in \mathscr{F} \implies P(E^{\text{C}}) = 1 - P(E)$ (\emph{probabilità del complementare});
                \item $E \in \mathscr{F} \implies P(E) \leq 1$;
                \item $E,F \in \mathscr{F} \land F \subset E \implies P(E \backslash F) = P(E) - P(F)$;
                \item $E,F \in \mathscr{F} \land F \subset E \implies P(F) \leq P(E)$ (\emph{monotonia});
                \item $E,F \in \mathscr{F} \implies P(E \cup F) = P(E) + P(F) - P(E \cap F)$ (\emph{probabilità dell'unione}).
            \end{enumerate}
        \end{prty}
        \begin{proof}
            \hfill
            \begin{enumerate}
                \item Osserviamo che, per l'assioma $(2)$ della Definizione~\ref{defn:Probabilità_impostazione_assiomatica}, vale $\Omega = E \cup E^{\text{C}}$ e sapendo che un evento e il suo complementare sono disgiunti: \[
                    1 = P(\Omega)= P(E) + P(E^{\text{C}}) \implies P(E^{\text{C}}) = 1 - P(E)
                .\] 
                \item Dalla precedente considerazione sappiamo che $P(E) = 1 - P(E^{\text{C}})$; aggiungendo l'assioma $(1)$ della Definizione~\ref{defn:Probabilità_impostazione_assiomatica} ($P(E^{\text{C}}) \geq 0$) segue che deve valere $P(E) \leq 1$.
                \item Vale $F \subset E \implies E = (E \backslash F) \cup F$ da cui deduciamo, per l'assioma $(2)$ della Proprietà~\ref{prty:Spazio_di_probabilità}, la seguente scrittura: \[
                    P(E) = P(E \backslash F) + P(F) \implies P(E \backslash F) = P(E) - P(F)
                .\] 
                \item Dalla precedente ricaviamo immediatamente $P(E) = P(E \backslash F) + P(F)$ e  $P(F) = P(E) - P(E \backslash F)$; sapendo dall'assioma $(1)$ della Definizione~\ref{defn:Probabilità_impostazione_assiomatica} che tutti e tre i termini dell'uguaglianza sono positivi o al più nulli, abbiamo necessariamente $P(F) \leq P(E)$.
                \item Scriviamo l'unione disgiunta dei due eventi $E$ ed  $F$ nel modo seguente:  \[
                    E \cup F = (E \cap F^{\text{C}}) \cup (E \cap F) \cup (E^{\text{C}} \cap F)
                .\] Per il punto $(2)$ della Proprietà~\ref{prty:Spazio_di_probabilità} la probabilità dell'unione vale: \[
                    P(E \cup F) = P(E \cap F^{\text{C}}) + P(E \cap F) + P(E^{\text{C}} \cap F)
                .\] Riscriviamo la precedente come:
                \begin{align*}
                    P(E \cup F) + P(E \cap F) &= \overset{P(E)}{\overbrace{P(E \cap F^{\text{C}}) + P(E \cup F)}} + \overset{P(F)}{\overbrace{P(E \cup F) + P(E^{\text{C}} \cap F)}} \\ {}&= P(E) + P(F)
                .
                \end{align*}
                Riordinando i termini, segue:
                \[
                    P(E \cup F) = P(E) + P(F) - P(E \cap F)
                .\] 
                Per ottenere le semplificazioni evidenziate con le graffe abbiamo usato le proprietà delle operazioni tra eventi (§~\ref{sec:Eventi}). \qedhere
            \end{enumerate}
        \end{proof}
    \section{Spazi finiti e numerabili}
    \begin{prty}\label{prty:Spazi_numerabili}
        Sia $\Omega$ uno spazio campionario \underline{numerabile} e  $\{\omega_1,\,\omega_2,\,\ldots\}$ una numera\-zione dei suoi punti (eventi elementari); scegliamo come $\sigma$\nbdash algebra  $\mathscr{F}$ l'insieme di tutti i sottoinsiemi (\emph{insieme delle parti}) dello spazio campionario  $\mathscr{P}(\Omega)$; si ha che:
            \begin{enumerate}
                \item Ogni probabilità $P({\omega_k})$ di un evento $E_k = \{\omega_k\} \in \mathscr{F}$ su uno spazio probabiliz\-zabile $(\Omega,\,\mathscr{F})$ individua una successione di numeri reali $p_1,\,p_2,\,\ldots$ che soddisfano la Proprietà~\ref{prty:Proprietà_funzione_probabilità} se scriviamo che $\forall k \in [1,\,\infty)\,\:\,P({\omega_k}) = p_k$;
                \item Data una successione $p_1,\,p_2,\,\ldots$ che soddisfa la Proprietà~\ref{prty:Proprietà_funzione_probabilità} esiste un'unica probabilità su $(\Omega,\,\mathscr{F})$ tale che  $P({\omega_k}) = p_k$ per ogni  $k$; essa è data da: \[
                        \forall E \subset \Omega\,:\,P(E) = \sum_{k\,:\,\omega_k \in E} p_k
                .\] 
            \end{enumerate}
        \end{prty}
        \begin{obsv}
            Nel caso di $\Omega$ spazio campionario \underline{finito} e numerabile, la Proprietà~\ref{prty:Spazi_numerabili} continua a valere, rispetto alla cardinalità di $\Omega$.
        \end{obsv}
    \section{Probabilità condizionata}
    \begin{defn}\label{defn:Probabilità_condizionata}
            Sia dato uno spazio di probabilità $(\Omega,\,\mathscr{F},\,P)$ e un evento $F \in \mathscr{F}$ tale che $P(F) > 0$; preso un qualsiasi altro evento $E \in \mathscr{F}$, si chiama \emph{probabilità condizionata} dell'evento $E$ dato il verificarsi di $F$:
            \begin{equation}\label{eq:Formula_probabilità_condizionata}
                P(E|F) = \frac{P(E \cap F)}{P(F)}
            .
            \end{equation}
        \end{defn}
    \section{Formula delle probabilità totali}
        \begin{defn}\label{defn:Probabilità_totali}
            Consideriamo uno spazio di probabilità $(\Omega,\,\mathscr{F},\,P)$ e una partizione finita $F_1,\,\ldots,\,F_n \in \mathscr{F}$ di $\Omega$; valga inoltre $\bigcup_{k=1}^{n} F_k = \Omega$ e infine $\forall h,k \in [1,\,n] \,:\, F_h \cap F_k = \text{\O} \land P(F_k) > 0$.
            Allora per qualunque evento $E \in \mathscr{F}$ la sua probabilità è definita come:
            \begin{equation}\label{eq:Formula_probabilità_totali}
                P(E) = \sum_{k=1}^{n} P(E|F_k) \cdot P(F_k)
            . 
            \end{equation}
        \end{defn}
        \begin{proof}
            Prendiamo l'evento $E \in \mathscr{F}$; dalle considerazioni fatte sopra abbiamo la seguente implicazione:  \[
                \Omega = \bigcup_{k=1}^{n} F_k \land E \subset \Omega \implies E = E \cap \Omega = \bigcup_{k=1}^{n} (E \cap F_k)
            .\] 
            Inoltre, poiché le partizioni $F_1,\,\ldots,\,F_n$ sono \emph{disgiunte a coppie}, la precedente $\bigcup_{k=1}^{n} (E \cap F_k)$ è un'unione disgiunta, e applicando il punto $(2)$ della Proprietà~\ref{prty:Spazio_di_probabilità} ($\bigstar$) e l'Equazione~\ref{eq:Formula_probabilità_condizionata} ($\ast$) otteniamo: \[
                P(E) = P\left(\bigcup_{k=1}^{n} (E \cap F_k)\right) \overset{\bigstar}{=} \sum_{k=1}^{n} P(E \cap F_k) \overset{\ast}{=} \sum_{k=1}^{n} P(E|F_k) \cdot P(F_k)
            .\] 
            Abbiamo così ottenuto la formula \eqref{eq:Formula_probabilità_totali}.
        \end{proof}
        \begin{obsv}
            La formula delle probabilità totali è utile quando le condizioni di preparazione di un esperimento aleatorio sono anch'esse casuali, e determinano una partizione dello spazio di probabilità dell'esperimento.
        \end{obsv}
    \section{Formula di Bayes}
        \begin{defn}\label{defn:Formula_Bayes}
            Consideriamo uno spazio di probabilità $(\Omega,\,\mathscr{F},\,P)$ e una partizione finita $F_1,\,\ldots,\,F_n \in \mathscr{F}$ di $\Omega$ tale che $\forall k \in [1,\,n] \,:\, P(F_k) > 0$; se abbiamo un evento $E \in \mathscr{F}$ per il quale $P(E) > 0$ allora otteniamo:
            \begin{align}\label{eq:Formula_Bayes}
                P(F_h|E) &= \frac{P(E|F_h) \cdot P(F_h)}{\sum_{k=1}^{n} P(E|F_k) \cdot P(F_k)} & h &= 1,\,\ldots,\,n
            .
            \end{align}
        \end{defn}
        \begin{proof}
            Usando \eqref{eq:Formula_probabilità_condizionata} possiamo scrivere: \[
                P(F_h|E) = \frac{P(F_h \cap E)}{P(E)} = \frac{P(E|F_h) \cdot P(F_h)}{P(E)}
            .\] 
            Applicando al denominatore \eqref{eq:Formula_probabilità_totali} otteniamo proprio la scrittura \eqref{eq:Formula_Bayes}.
        \end{proof}
        \begin{obsv}
            La formula di Bayes è utile quando possediamo informazioni a posteriori su un esperimento aleatorio e vogliamo determinare le condizioni entro le quali si sia verificato un certo evento.
        \end{obsv}
    \section{Formula di moltiplicazione}
        \begin{defn}\label{defn:Formula_moltiplicazione}
            Consideriamo uno spazio di probabilità $(\Omega,\,\mathscr{F},\,P)$ e una successione di eventi al suo interno $E_1,\,\ldots,\,E_n \in \mathscr{F}$, per i quali valga $P(\bigcap_{k=1}^{n-1} E_k) > 0$; allora possiamo scrivere:
            \begin{equation}\label{eq:Formula_moltiplicazione}
                P(E_1 \cap \dotsm \cap E_n) = P(E_1) \cdot P(E_2|E_1) \cdot P(E_3|E_2 \cap E_1) \cdot \ldots \cdot P(E_n|E_1 \cap \dotsm \cap E_{n-1})
            .
            \end{equation}
        \end{defn}
        \begin{proof}
            Dato che $(E_1 \cap \dotsm \cap E_{n-1}) \subset (E_1 \cap \dotsm \cap E_{n-2}) \subset \dotsm \subset E_1$ usando il punto $(4)$ della Proprietà~\ref{prty:Proprietà_funzione_probabilità} vale: \[
                0 < P(E_1 \cap \dotsm \cap E_{n-1}) \leq P(E_1 \cap \dotsm \cap E_{n-2}) \leq \dotsm \leq P(E_1)
            .\]
            Otteniamo quindi il seguente prodotto:
            \begin{align*}
                P(E_1 \cap \dotsm \cap E_n) &= P(E_1) \cdot \frac{P(E_1 \cap E_2)}{P(E_1)} \cdot \frac{P(E_1 \cap E_2 \cap E_3)}{P(E_1 \cap E_2)} \cdot \ldots \cdot \frac{P(E_1 \cap \dotsm \cap E_n)}{P(E_1 \cap \dotsm \cap E_{n-1})} \\ &= P(E_1) \cdot P(E_2|E_1) \cdot P(E_3|E_2 \cap E_1) \cdot \ldots \cdot P(E_n|E_1 \cap \dotsm \cap E_{n-1})
            .
            \end{align*}
            Si noti che abbiamo usato \eqref{eq:Formula_probabilità_condizionata} per riscrivere il prodotto precedente, ottenendo \eqref{eq:Formula_moltiplicazione}.
        \end{proof}
    \section{Eventi indipendenti}
        \begin{defn}\label{defn:Eventi_indipendenti}
            Sia $(\Omega,\,\mathscr{F},\,P)$ uno spazio di probabilità; gli eventi $E,\,F \in \mathscr{F}$ sono \emph{indipendenti} se vale: \[
                P(E \cup F) = P(E) \cdot P(F)
            .\] 
        \end{defn}
        \begin{obsv}
            Se due eventi $E$ ed $F$, presi dallo stesso spazio di probabilità, sono indipendenti, allora valgono le seguenti uguaglianze:
            \begin{itemize}
                \item $P(E|F) = P(E)$;
                \item $P(F|E) = P(F)$.
            \end{itemize}
        \end{obsv}
        \begin{defn}
            Sia $(\Omega,\,\mathscr{F},\,P)$ uno spazio di probabilità; diciamo che gli eventi $E_1,\, \ldots,\, E_n$ sono \emph{indipendenti} se comunque preso un sottoinsieme $\{h_1,\, \ldots,\, h_k\} \subset \{1,\, \ldots,\, n\}$ con $k \geq 2$ vale la seguente uguaglianza: \[
                P(E_{h_1} \cap \dotsm \cap E_{h_k}) = P(E_{h_1}) \cdot \ldots \cdot P(E_{h_k})
            .\] 
        \end{defn}
        \begin{obsv}
            Per testare l'indipendenza di $m$ eventi, sarà necessario provare $2^{m} - m - 1$ uguaglianze (in dipendenza dalle possibili combinazioni di intersezioni tra gli eventi da testare).
        \end{obsv}
        \begin{defn}
            Preso uno spazio di probabilità e considerata una successione dei suoi eventi, diremo che essa è costituita da eventi indipendenti se, comunque scelto un sottoinsieme finito di eventi dalla successione, esso è costituito da eventi indipendenti.
        \end{defn}
        \begin{defn}
            Consideriamo uno spazio di probabilità $(\Omega,\,\mathscr{F},\,P)$, una successione di eventi $A_1,\, \ldots,\, A_n$ e un evento $F$ tale che $P(F) > 0$; gli eventi della successione si dicono \emph{condizionatamente indipendenti} dato $F$ se essi sono indipendenti rispetto alla probabilità $P(F)$.
        \end{defn}
    \section{Affidabilità dei sistemi}
        \begin{defn}
            La probabilità che un componente di un sistema non si guasti durante il periodo di tempo in cui deve operare è detta \emph{affidabilità} del componente; essa viene espressa rispetto all'interazione che i componenti hanno tra loro:
            \begin{itemize}
                \item \textbf{serie}: per garantire il funzionamento del sistema tutti i componenti collegati in serie devono funzionare correttamente; in questo caso vale \[
                    r = r_1 \cdot \ldots \cdot r_k
                ,\] per un sistema con $k$ componenti in serie;
            \item \textbf{parallelo}: per garantire il funzionamento del sistema basta che almeno un componente funzioni correttamente; in questo caso vale \[
                        r = 1 - (1 - r_1) \cdot \ldots \cdot (1 - r_k)
                ,\] per un sistema con $k$ componenti in parallelo.
            \end{itemize}
            In entrambi i casi appena mostrati, $r_i$ rappresenta l'affidabilità dell'$i$\nbdash esimo componente del sistema.
        \end{defn}
        \begin{obsv}
            Per analizzare problemi sull'affidabilità è conveniente scomporre un sistema complesso in sottosistemi che abbiano solo componenti in serie o solo in parallelo.
        \end{obsv}
    \section{Prove di Bernoulli}
        \begin{defn}
            Siano $n \in \mathbb{N}$ e $p \in (0,\,1)$; consideriamo il seguente spazio campionario: \[
                \Omega \coloneqq \{(a_1,\,\ldots,\,a_n)\,:\,a_k \in \{0,\,1\},\, k \in [1,\,n]\}
                ,\] la $\sigma$\nbdash algebra $\mathscr{F} = \mathscr{P}(\Omega)$ e la funzione di probabilità: \[
                \forall (a_1,\, \ldots,\, a_n) \in \Omega \,:\, P(\{a_1,\, \ldots,\, a_n\}) = p^{\left[\sum_{k=1}^{n} a_k\right]} \cdot (1 - p)^{\left[n - \sum_{k=1}^{n} a_k\right]}
            ;\] 
            la terna $(\Omega,\,\mathscr{F},\,P)$ si chiama \emph{spazio di probabilità di Bernoulli} oppure spazio di probabilità di $n$ prove di Bernoulli.
        \end{defn}
        \begin{prty}
            Consideriamo uno spazio di probabilità di Bernoulli nel quale la probabilità di successo della singola prova è $p \in (0,\,1)$; la probabilità di osservare $k \leq n$ successi in una sequenza di $n \geq 1$ prove di Bernoulli in questo spazio è data da:
            \begin{align}\label{eq:Probabilità_Bayes}
                \binom{n}{k} \cdot p^{k} \cdot (1 - p)^{n-k}
            .
            \end{align}
        \end{prty}
        \begin{proof}
            Sia $(\Omega,\,\mathscr{F},\,P )$ uno spazio di probabilità di $n$ prove Bernoulli e $B_k \in \mathscr{F}$ l'evento `osserviamo $k$ successi in $n$ prove', ovvero: \[
                B_k = \left\{ (a_1,\, \ldots,\, a_n) \in \Omega \,:\, \sum_{h=1}^{n} a_h = k \right\} 
            ;\] la probabilità di questo evento è definita come: \[
                P(B_k) = \sum_{\omega \in B_k} P(\{\omega\}) = \sum_{\omega \in B_k} p^{k}(1 - p)^{n - k} = |B_k|p^{k}(1 - p)^{n - k}
            .\] Osservando che $|B_k| = \binom{n}{k}$ (tutte le stringhe di $n$ bit contenenti $k$ zeri e $n - k$ uni) otteniamo l'equazione \eqref{eq:Probabilità_Bayes}
        \end{proof}
    \begin{obsv}
        Gli eventi $B_k$ che fissano $k$ successi in $n$ prove di Bernoulli hanno delle probabilità che corrispondono al modello binomiale $p_k$; possiamo affermare che, preso lo spazio campionario $\hat{\Omega} = \{0,\,\ldots,\,n\}$ dell'esperimento che considera $k$ successi in $n$ prove, lo spazio di probabilità di Bernoulli induce su $\hat{\Omega}$ un modello binomiale di parametri $n$ e $p$.
    \end{obsv} 
    \section{Serie geometrica}
    \begin{defn}[Serie]
        Data una successione $\{a_{n}\}$, se sommiamo gli infiniti termini otteniamo una \emph{serie}, definita come $\sum_{n=0}^{\infty} a_n$ (serie di a con n da 0 a infinito).

        Si dice somma parziale $n$\nbdash esima di una serie $s_{n}=\sum_{n=0}^{\infty} a_i,\,n \in \mathbb{N}$ (la serie $\sum_{n=0}^{\infty} a_n$ converge, diverge o è irregolare se la successione delle sue somme parziali converge, diverge o è irregolare).

        Si dice somma di una serie il limite, se esiste finito, della successione delle sue somme parziali: $\lim_{n \to \infty} s_n = s$

        Condizione \underline{necessaria} (non sufficiente) affinché una serie converga è che il suo termine generale ($a_n$) tenda a 0
    \end{defn}
    \begin{defn}[Serie geometrica]\label{defn:Serie_geometrica}
        Chiamiamo \emph{geometrica} la serie col seguente termine generale:
        \begin{align*}
            \sum_{i=0}^{n} q^{i} \cdot \frac{1 - q^{n+1}}{1 - q} & &q\neq 1
        .\end{align*}
        La somma della serie con ragione $q \in (0,\,1)$ vale: \[
            \sum_{i=0}^{n} q^{i} = \begin{cases}
                \frac{1}{1-q} & \text{se $|q| < 1$}; \\
                +\infty & \text{se $q > 1$}; \\
                \text{indeterminata} & \text{altrimenti}
            \end{cases}
        .\] 
    \end{defn}

    %! TEX root = main.tex
% Capitolo 2

\chapter{Variabili Aleatorie}

    %! TEX root = main.tex
% Capitolo 3

\chapter{Vettori Aleatori}
    \section{Vettore aleatorio}
        \begin{defn}
            Consideriamo $n$ variabili aleatorie $X_1, \ldots,\, X_n$ definite tutte sullo stesso spazio di probabilità $(\Omega,\,\mathbb{F},\,P)$; diremo che sono  \emph{indipendenti} se vale:
            \begin{equation}\label{eq:Indipendenza_variabili_aleatorie}
                P(X_1 \in B_1,\, \ldots,\, X_n \in B_n) = P(X_1 \in B_1) \cdot \ldots \cdot P(X_n \in B_n)
            ,\end{equation}
            per ogni scelta dei domini regolari $B_1,\, \ldots,\, B_n$ ottenuti al più con un numero finito e numerabile di operazioni tra intervalli.
        \end{defn}
        \begin{obsv}
            Se prendiamo dei domini regolari definiti come: \[
                \forall i \in (-\infty,\, \ldots x_i]
            ,\] l'equazione \eqref{eq:Indipendenza_variabili_aleatorie} diventa:
            \begin{equation}\label{eq:Indipendenza_v_a_ripartizione}
                P(X_1 \leq x\,\, \ldots,\, X_N \leq x_n) = P(X_1 \leq x_1) \cdot \ldots \cdot P(X_n \leq x_n)
            .\end{equation}
        \end{obsv}
        \begin{prty}
            Prese $n$ variabili aleatorie $X_1,\, \ldots,\, X_n$, esse sono indipendenti se vale  \eqref{eq:Indipendenza_v_a_ripartizione} per ogni scelta di $x_1,\, \ldots,\, x_n \in \mathbb{R}$.
        \end{prty}
        \begin{obsv}
            Consideriamo $n$ variabili aleatorie indipendenti e discrete $X_1,\, \ldots,\, X_n$ con densità rispettivamente $p_{X_1},\, \ldots,\, p_{X_n}$; allora scegliendo i domini in \eqref{eq:Indipendenza_variabili_aleatorie} come: \[
                B_1 = \{x_1\},\, \ldots,\, B_n = \{x_n\}
            ,\] otteniamo che:
            \begin{align}\label{eq:Indipendenza_v_a_densità}
                P(X\ = x_1,\, \ldots,\, X_n = x_n) = P(X_1 = x_1) \cdot \ldots \cdot P(X_n = x_n) & &
                \forall  x_i \in \mathbb{R} \forall i = 1,\, \ldots,\, n
            .\end{align}
        \end{obsv}
        \begin{prty}
            Possiamo affermare che, prese $n$ variabili aleatorie discrete $X_1,\, \ldots,\, X_n$, esse sono indipendenti se vale \eqref{eq:Indipendenza_v_a_densità}.
        \end{prty}
    \section{Vettori aleatori}
        \begin{defn}
            Sia dato uno spazio di probabilità $(\Omega,\,\mathbb{F},\,P)$; un \emph{vettore aleatorio} $n$\nbdash dimensionale è una funzione vettoriale definita come:
            \begin{align*}
                \vec{X} \coloneqq (X_1,\, \ldots,\, X_{n}) & &
                \vec{X}\,:\, \Omega \mapsto \mathbb{R}^n
            ,\end{align*}
            tale che per ogni $i \in 1,\, \ldots,\, n$ ciascun $X_i$ sia una variabile aleatoria.
        \end{defn}
    \section{Funzione di ripartizione congiunta}
        \begin{defn}
            Sia $\vec{X} = (X_1,\, \ldots,\, X_{n})$ un vettore aleatorio $n$\nbdash dimensionale definito sullo spazio di probabilità $(\Omega,\,\mathbb{F},\,P)$; chiamiamo funzione di ripartizione di $\vec{X}$ (oppure funzione di ripartizione \emph{congiunta} di $X_1,\, \ldots,\, X_{n}$) la funzione: \[
                F_{\vec{X}} = F_{(X_1,\, \ldots,\, X_{n})} \,:\, \mathbb{R}^n \mapsto [0,\,1]
            ,\] definita $\forall x \in (x_1,\, \ldots,\, x_{n}) \in \mathbb{R}^n$ come: \[
                F_{(X_1,\, \ldots,\, X_{n})}(x_1,\, \ldots,\, x_{n}) \coloneqq P(X_1 \leq x_1,\, \ldots,\, X_n \leq x_n)
            .\]
        \end{defn}
        \begin{prty}\label{prty:Convergenza_ripartizione_congiunta}
            Sia $\vec{X} = (X_1,\, \ldots,\, X_{n})$ un vettore aleatorio che ammette funzione di ripartizione $F_{\vec{X}}$ e sia $\vec{x} = (x_1,\, \ldots,\, x_{n})$; allora vale: \[
                \forall k \in [1,\,n] \,:\, \lim_{x_k \to \infty} F_{\vec{X}}(\vec{x}) = 0
            ,\] mentre otteniamo:
            \begin{align*}
                \lim_{x_k \to \infty} F_{\vec{X}}(\vec{x}) &= 
                P(X_1 \leq x_1,\, \ldots,\, X_{k-1} \leq x_{k-1},\, X_{k+1} \leq x_{k+1},\, \ldots,\, X_n \leq x_n) \\
                                                           &= F_{(X_1,\, \ldots,\, X_{k-1},\, X_{k+1},\, \ldots,\, X_{n})}(x_1,\, \ldots,\, x_{k-1},\, x_{k+1},\, \ldots,\, x_{n})
            .\end{align*}
        \end{prty}
    \section{Indipendenza di variabili aleatorie}
        \begin{obsv}
            Consideriamo un vettore aleatorio bidimensionale $(X,\,Y)$ con funzione di ripartizione $F_{X,\,Y}$; la Proprietà~\ref{prty:Convergenza_ripartizione_congiunta} afferma che:
            \begin{gather*}
                \lim_{x \to \infty} F_{X,\,Y}(x,\,y) = P(Y \leq y) = F_Y(y); \\
                \lim_{y \to \infty} F_{X,\,Y}(x,\,y) = P(X \leq x) = F_X(x)
            .\end{gather*}
            Se prendiamo in generale un vettore aleatorio $n$\nbdash dimensionale, applicando la Proprietà~\ref{prty:Convergenza_ripartizione_congiunta} iterativamente, otteniamo:
            \begin{align*}
                F_{X_i}(x) = \lim_{\substack{x_j \rightarrow +\infty \\ \forall j \neq i}} F_{\vec{X}}(x_1,\, \ldots,\, x_{i-1},\, x,\, x_{i+1},\, \ldots,\, x_{n})
            .\end{align*}
            Concludiamo che dalla funzione di ripartizione congiunta si possono calcolare le ripartizioni marginali, ma non vale il contrario.
        \end{obsv}
        \begin{defn}
            Le componenti di un vettore aleatorio $\vec{X} = (X_1,\, \ldots,\, X_{n})$ sono indipendenti se e solo se la funzione di ripartizione di $\vec{X}$ coincide col prodotto delle ripartizioni marginali: \[
            F_{\vec{X}} = F_{X_1} \cdot \ldots \cdot F_{X_n}
            .\] 
        \end{defn}
    \section{Vettori aleatori discreti}
        \begin{defn}
            Un vettore aleatorio $n$\nbdash dimensionale $X$ è \emph{discreto} se le sue componenti $X_1,\, \ldots,\, X_{n}$ sono variabili aleatorie discrete.
        \end{defn}
    \section{Densità discrete}
        \begin{defn}\label{defn:Densità_congiunta}
            Sia $\vec{X}$ un vettore aleatorio discreto su uno spazio di probabilità $(\Omega,\,\mathbb{F},\,P)$; la funzione: \[
                p_{\vec{X}} \coloneqq P(X_1 = x_1,\, \ldots,\, X_n = x_n)
            ,\] dove vale $\vec{x} = (x_1,\, \ldots,\, x_{n})$, si chiama  \emph{densità discreta} del vettore aleatorio $\vec{X}$ oppure densità congiunta di $X_1,\, \ldots,\, X_{n}$.
        \end{defn}
    \section{Distribuzione multinomiale}
        \begin{defn}\label{defn:Densità_multinomiale}
            Consideriamo una popolazione che contenga $k \geq 2$ tipi di oggetti, dove la proporzione degli oggetti di tipo  $i$\nbdash esimo sul totale è rappresentata da $p_i$ come:  \[
                \forall i \in [1,\, k] \,:\, p_i > 0 \land \sum_{i=1}^{k} p_i = 1
            ;\] inoltre $n$ oggetti di questa popolazione siano estratti a caso con reimmissione.
            Sia $X_i$ il numero di oggetti di tipo $i$\nbdash esimo estratti (con $i \in [1,\, k]$) e sia $\vec{X}$ il vettore aleatorio con componenti $X_1,\, \ldots,\, X_{k}$; allora il vettore $\vec{X}$ è discreto e tale che: \[
            X_1 + \ldots + X_{k} = n
            ,\] e la sua densità è detta \emph{multinomiale} di parametri $n$ e $p_1,\, \ldots,\, p_{k}$:
            \begin{align}\label{eq:Densità_multinomiale}
                P(X_1 = n_1,\, \ldots,\, X_{k} = n_k) &= \binom{n}{n_1,\, \ldots,\, n_{k}} \cdot (p_1^{n_1} \cdot \ldots \cdot p^{n_k}_{k}) \\
                                                      & \text{per } n_1,\, \ldots,\, n_{k} \in [0,\, n] \land n_1 + \ldots + n_k = n \nonumber
            .\end{align}
        \end{defn}
        \begin{prty}\label{prty:Densità_congiunta}
            Sia $p_{\vec{X}}$ la densità di un vettore aleatorio $n$\nbdash dimensionale $\vec{X}$ che assume valori in un insieme al più numerabile $S$ con probabilità 1; allora possiamo dire che:
            \begin{enumerate}
                \item $\forall \vec{x} \in \mathbb{R}^n \,:\, 0 \leq p_{\vec{X}} \leq 1 \land 
                    \forall \vec{x} \notin S \,:\, p_{\vec{X}}(\vec{x}) = 0$;
                \item $\sum_{\vec{x} \in S} è_{\vec{X}}(\vec{x}) = 1$;
                \item se $F_{\vec{X}}$ è la funzione di ripartizione di $\vec{X}$, allora: \[
                    \forall \vec{x} \in \mathbb{R}^n \,:\, F_{\vec{X}}(\vec{x}) = \!\!\!\! \sum_{\vec{y} \in S \,:\, \vec{y} \leq \vec{x}} \!\!\!\! p_{\vec{X}}(\vec{y})
                ;\] 
                \item se $B \subset \mathbb{R}^n$, allora: \[
                        P(\vec{X} \in B) = \!\!\! \sum_{\vec{x} \in B \cap S} \!\!\! p_{\vec{X}}(\vec{x})
                .\] 
            \end{enumerate}
        \end{prty}
        \begin{proof}
            La dimostrazione segue lo stesso svolgimento di quella della Proprietà~\ref{prty:Densità_discreta}.
        \end{proof}
        \begin{defn}\label{defn:Densità_marginali}
            Chiamiamo \emph{densità marginali} le densità delle componenti $X_i$ del vettore aleatori $\vec{X}$; se prendiamo la prima componente $X_1$, ricordando che:
            \begin{align*}
                p_{X_1}(x_1) &= P(X_1 = x_1) = P(X_1 = x_1,\, X_2 \in \mathbb{R},\, \ldots,\, X_n \in \mathbb{R}) \\
                             &= P(\vec{X} \in B)
            ,\end{align*}
            dove $B \coloneqq \{x_1\} \times \mathbb{R}^{n-1}$; quindi possiamo scrivere la densità come:
            \begin{equation}\label{eq:Densità_marginali}
                p_{X_1}(x_1) = \!\! \sum_{\vec{x} \in B \cap S} \!\! p_{\vec{X}}(\vec{x}) 
                = \!\! \sum_{x_2,\, \ldots,\, x_{n}} \!\! p_{\vec{X}}(x_1,\, x_2,\, \ldots,\, x_{n})
            .\end{equation}
        \end{defn}
        \begin{prty}\label{prty:Indipendenza_componenti_densità}
            Le componenti di un vettore aleatorio discreto $\vec{X} = (X_1,\, \ldots,\, X_{n})$ sono indipendenti se e solo se la densità congiunta di $\vec{X}$ coincide col prodotto di delle densità marginali $p_{X_1},\, \ldots,\, p_{X_n}$ rispettivamente di $X_1,\, \ldots,\, X_{n}$, cioè: \[
            p_{\vec{X}} = p_{X_1} \cdot \ldots \cdot p_{X_n}
            .\]
        \end{prty}
    \section{Vettori aleatori assolutamente continui}
        \begin{defn}\label{defn:Vettori_aleatori_assolutamente_continui}
            Un vettore aleatorio $\vec{X}$ $n$\nbdash dimensionale è \emph{assolutamente continuo} se esiste una funzione $f_{\vec{X}}\,:\, \mathbb{R}^n \mapsto \mathbb{R}^+$ integrabile, tale che la funzione di ripartizione di $\vec{X}$ si possa scrivere come:
            \begin{align*}
                F_{\vec{X}}(\vec{x}) = \int_{-\infty}^{x_1} \dotsi 
                \int_{-\infty}^{x_n} f_{\vec{X}}(S_1,\, \ldots,\, S_{n})\, d_{S_n} \cdot \ldots \cdot d_{S_1} 
                & & \forall \vec{x} = (x_1,\, \ldots,\, x_{n})
            .\end{align*}
            La funzione $f_{\vec{X}}$ prende il nome di \emph{densità congiunta} di $X_1,\, \ldots,\, X_{n}$ oppure densità del vettore aleatorio assolutamente continuo $\vec{X}$.
        \end{defn}
        \begin{prty}\label{prty:Vettori_aleatori_assolutamente_continui}
            Sia $f_{\vec{X}}$ la densità di un vettore aleatorio $n$\nbdash dimensionale assolutamente continuo; allora possiamo affermare che:
            \begin{enumerate}
                \item $\int_{\mathbb{R}^n} f_{\vec{X}}(x_1,\, \ldots,\, x_{n})\, dx_1 \cdot \ldots \cdot dx_n = 1$;
                \item se  $F_{\vec{X}}$ è la funzione di ripartizione di $\vec{X}$, allora: \[
                    \frac{\partial^n F_{\vec{X}}(\vec{x})}{\partial x_1 \cdot \ldots \cdot \partial x_{n}} = f_{\vec{X}}(\vec{x})
                ,\] per ogni $\vec{x} \in \mathbb{R}^n$ tale che esista la derivata parziale al primo membro;
                \item se $B \subset \mathbb{R}^n$ è un \underline{dominio regolare}, allora: \[
                        P(\vec{X} \in B) = \int_{B} f_{\vec{X}}(x_1,\, \ldots,\, x_{n})\, dx_1 \cdot \ldots \cdot dx_n
                .\] 
            \end{enumerate}
        \end{prty}
        \begin{proof}
            La dimostrazione segue lo stesso svolgimento di quella della Proprietà~\ref{prty:Variabile_aleatoria_continua}.
        \end{proof}
        \begin{prty}[Densità marginale continua]\label{prty:Densità_marginale_continua}
            Se $f_{\vec{X}}$ è la densità di un vettore aleatorio $n$\nbdash dimensionale assolutamente continuo $\vec{X} = (X_1,\, \ldots,\, X_{n})$ allora $X_i$ è una variabile aleatoria assolutamente continua e la sua densità si ottiene come: \[
                f_{X_i}(x_i) = \int_{\mathbb{R}^{n-1}} f_{\vec{X}}(s_1,\, \ldots,\, s_{i-1},\, x_i,\, s_{i+1},\, \ldots,\, s_n)\, ds_1 \cdot \ldots \cdot ds_{i-1} \cdot ds_{i+1} \ldots \cdot ds_n
            .\] 
        \end{prty}
        \begin{proof}
            Consideriamo il caso $i=1$; dobbiamo dimostrare che: \[
            F_{X_i}(x) = \int_{-\infty}^{\infty} \left\{\int_{\mathbb{R}^{n-1}} f_{\vec{X}}(s_1,\, \ldots,\, s_n) ds_2 \cdot \ldots \cdot ds_n\right\}\, ds_1
            ;\] ciò si verifica dato che, se definiamo $B \coloneqq (-\infty,\, x] \times \mathbb{R}^{n-1}$, allora per il punto $(3)$ della Proprietà~\ref{prty:Vettori_aleatori_assolutamente_continui} otteniamo:
            \begin{align*}
                &\int_{-\infty}^{\infty} \left\{\int_{\mathbb{R}^{n-1}} f_{\vec{X}}(s_1,\, \ldots,\, s_n) ds_2 \cdot \ldots \cdot ds_n\right\}\, ds_1 = P(\vec{X} \in B) \\
                &= P(X_1 \leq x_1,\, X_2 \in \mathbb{R},\, \ldots,\, X_n) = P(X_1 \leq x) = F_{X_1}(x) \qedhere
            .\end{align*}
        \end{proof}
        \begin{prty}
            Le componenti di un vettore aleatorio assolutamente continuo sono indipendenti se e solo se ammettono densità congiunta che può essere ottenuta come prodotto delle densità marginali.
        \end{prty}
    \section{Funzioni di vettori aleatori discreti}
    \begin{defn}\label{defn:Funzione_discreta}
            Sia $\vec{X}$ un vettore aleatorio discreto con densità $p_{\vec{X}}(\vec{x})$ per cui valga $P(\vec{X} \in S) = 1$ con $S$ al più numerabile; consideriamo inoltre $\vec{g}\,:\, S \mapsto \mathbb{R}^m$ e $\vec{Y} \coloneqq \vec{g}(\vec{X})$.

            $Y$ è un vettore aleatorio discreto, per i cui valori otteniamo:  \[
                g(S) = \left\{\vec{y} = \vec{g}(\vec{x}) \,:\, \vec{x} \in S\right\} \implies 
                P(\vec{Y} \in \vec{g}(S)) = 1
            .\]
        \end{defn}
        \begin{prty}[Densità di funzione discreta]\label{prty:Densità_funzione_discreta}
            Sia $\vec{X}$ un vettore aleatorio discreto con densità $p_{\vec{X}}(\vec{x})$ per cui valga $P(\vec{X} \in S) = 1$ con $S$ al più numerabile, e consideriamo $\vec{g}\,:\, S \mapsto \mathbb{R}^m$; allora $\vec{Y} \coloneqq  \vec{g}(\vec{X})$ è un vettore aleatorio discreto tale che $P(\vec{Y} \in \vec{g}(S)) = 1$ e la densità di $\vec{Y}$ si ottiene come:
            \begin{equation}\label{eq:Densità_funzione_discreta}
                p_{\vec{y}} = \!\!\! \sum_{\substack{\vec{x} \in S \,:\\ 
                \vec{g}(\vec{x}) = \vec{y}}} \!\! p_{\vec{X}}(\vec{x})
            .\end{equation}
        \end{prty}
        \begin{prty}[Somma discreta]\label{prty:Somma_variabili_aleatorio_discrete}
            Sia $(X_1,\, X_2)$ un vettore aleatorio discreto con densità $p_{X_1,\,X_2}(x_1,\, x_2)$ e prendiamo $Y\coloneqq X_1 + X_2$; dalla formula \eqref{eq:Densità_funzione_discreta} otteniamo che: \[
                p_{Y} = p_{X_1 + X_2}(y) = \!\!\!\! \sum_{\substack{x_1,\,x_2 \,:\\
                x_1 + x_2 = y}} \!\!\! p_{X_1,\, X_2}(x_1,\, x_2) = 
                \sum_{x_2}p_{X_1,\, X_2}(y - x_2,\, x_2)
            .\] In particolare, se $X_1,\, X_2$ sono \underline{indipendenti} allora possiamo scrivere: \[
            p_{X_1,\, X_2}(y) = \sum_{x_2} p_{X_1}(y-x_2)\cdot p_{X_2}(x_2)
            .\]  Possiamo estendere per ricorrenza questo risultato al caso di un vettore aleatori discreto $n$\nbdash dimensionale.
        \end{prty}
        \begin{obsv}[Somma di Poisson indipendenti]
            Siano $X_1,\, X_2$ variabili aleatorio indipendenti con densità di Poisson di parametri $\lambda_{1},\, \lambda_{2}$ rispettivamente; allora la loro somma è una variabile aleatoria discreta con densità di Poisson di parametro $\lambda = \lambda_1 + \lambda_2$:
            \begin{align*}
                P(X_1 + X_2 = k) &= \sum_{j=0}^{k} p_{X_1,\, X_2}(k-j,\, j) = 
                \sum_{j=0}^{k} p_{X_1}(k-j)\cdot p_{X_2}(j) \\
                                 &= \sum_{j=0}^{k} \left[e^{-\lambda_1} \frac{\lambda_1^{k-j}}{(k-j)!}\right]
                                 \cdot e^{-\lambda_2}\frac{\lambda_2^{j}}{j!} = 
                                 \frac{e^{-(\lambda_1 + \lambda_2)}}{k!}\cdot \sum_{j=0}^{k}\binom{k}{j}\lambda^{k-j}_1 \cdot \lambda^j_2 \\
                                 &= e^{-(\lambda_1 + \lambda_2)}\cdot \frac{(\lambda_1 + \lambda_2)^k}{k!}
            .\end{align*}
            Iterando questo procedimento, possiamo estenderlo nel caso di $X_1,\, \ldots,\, X_{n}$ variabile aleatorie indipendenti con densità $X_1\sim \mathcal{P}(\lambda_1),\, \ldots,\, X_{n}\sim \mathcal{P}(\lambda_n)$; 
            allora vale: \[
                X_1 + \ldots + X_{n} \sim \mathcal{P}(\lambda_1 + \ldots + \lambda_n)
            .\]
        \end{obsv}
        \begin{obsv}[Somma di Bernoulliane indipendenti]
            Siano $X_1,\, \ldots,\, X_{n}$ variabili aleatorio indipendenti con densità di Bernoulli, di parametro $p \in (0,\, 1)$; allora la loro somma avrà la seguente distribuzione: \[
                X_1 + \ldots + X_{n} \sim \mathcal{B}e(n,\, p)
            .\] 
        \end{obsv}
        \begin{obsv}[Somma di binomiali indipendenti]
            Siano $X_1,\, \ldots,\, X_{k}$ variabili aleatorio indipendenti con densità binomiali $X_1\sim \mathcal{B}i(n_1,\,p),\, \ldots X_n\sim \mathcal{B}i(n_k,\,p)$, con parametro $p \in (0,\, 1)$; allora la loro somma avrà la seguente distribuzione: \[
                X_1 + \ldots + X_{n} \sim \mathcal{B}i(n_1+ \ldots +n_k,\, p)
            .\] 
        \end{obsv}
        \begin{defn}[Trasformazioni affini di vettori discreti]
            Sia $\vec{A}$ una matrice $n \times n$ invertibile, $\vec{b}$ un vettore colonna di dimensione $n$ (valga $\vec{b} \in \mathbb{R}^n$), $\vec{X}$ un vettore aleatorio discreto $n$\nbdash dimensionale con densità $p_{\vec{X}}$; valga $\vec{Y} = \vec{A} \cdot \vec{X} + \vec{b}$, intendendo $\vec{X},\, \vec{Y}$ come vettori colonna.

            Otteniamo la densità di $\vec{Y}$ applicando la formula \eqref{eq:Densità_funzione_discreta} alla trasformazione biunivoca $\mathbb{R}^n \mapsto \mathbb{R}^n$ definita come $\vec{g}(\vec{x}) = \vec{A} \cdot \vec{x} + \vec{b}$ con inversa $\vec{g}^{\,-1}(\vec{y}) = \vec{A}^{\,-1}(\vec{y} - \vec{b})$.

            Poiché $\vec{g}$ è biunivoca, possiamo affermare che:
            \begin{equation}\label{eq:Trasformazioni_vettori_discreti}
                p_{\vec{Y}}(\vec{y}) = P(\vec{Y} = \vec{y}) = p_{\vec{X}}(\vec{A}^{\,-1}(\vec{y} - \vec{b}))
            .\end{equation}
        \end{defn}
    \section{Funzioni di vettori aleatori assolutamente continui}
        \begin{defn}
            Sia $X$ un vettore aleatorio $n$\nbdash dimensionale continuo con densità $f_{\vec{X}}$ e per $\vec{g}\,:\, \mathbb{R}^n \mapsto \mathbb{R}^m$ valga $\vec{Y} = \vec{g}(\vec{X})$.
            
            Non è sempre detto che in queste condizioni applicando una funzione qualunque $\vec{g}$ al vettore aleatorio $\vec{X}$, si ottenga un $\vec{Y}$ che sia ancora un vettore aleatorio; infatti $\vec{g}$ dovrà essere opportunamente regolare.
        \end{defn}
        \begin{prty}
            Consideriamo un vettore aleatorio $\vec{Y} = \vec{g}(\vec{X})$; la sua funzione di ripartizione si ottiene osservando che: \[
                F_{\vec{Y}}(\vec{y}) = P(\vec{Y} \leq \vec{y}) = P(\vec{X}) \leq \vec{y}
            .\] Applicando il punto $(3)$ della Proprietà~\ref{prty:Vettori_aleatori_assolutamente_continui} possiamo esprimere la $F_{\vec{Y}}(\vec{y})$ in funzione della densità di $\vec{X}$:
            \begin{align}\label{eq:Ripartizione_funzione_continua}
                F_{\vec{Y}}(\vec{y}) = P(\vec{X} \in A) = \int_{A} f_{\vec{X}}(\vec{x})\, d\vec{x} & &
                \text{per } A \coloneqq \{\vec{x} \,:\, \vec{g}(\vec{x}) \leq \vec{y}\}
            .\end{align}
        \end{prty}
        \begin{prty}[Somma continua]
            Sia $(X_1,\, X_2)$ un vettore aleatorio assolutamente continuo con densità $f_{X_1,\,X_2}$; allora la funzione di ripartizione di $X_1 + X_2$ si ottiene applicando l'equazione \eqref{eq:Ripartizione_funzione_continua} alla funzione $g(x_1,\, x_2)=x_1+x_2$:
            \begin{align*}
                F_{X_1+X_2}(y) &= \int_{\substack{\{x_1,\,x_2\,:\\
                x_1+x_2 \leq y\}}} f_{X_1,\,X_2}(x_1,\,x_2)\, dx_1\, dx_2 \\
                               &= \int_{-\infty}^{\infty} \int_{-\infty}^{y-x_1} f_{X_1,\,X_2}(x_1,\,x_2)\,dx_1\, dx_2 \\
                               &= \int_{-\infty}^{y} \left(\int_{-\infty}^{\infty} f_{X_1,\,X_2}(x_1,\,x_2 - x_1)\, dx_1\right)\, dx_2
            .\end{align*}
            Notiamo che $X_1+X_2$ è una variabile aleatoria assolutamente continua, e la sua densità si ottiene della ripartizione appena scritta: \[
                f_{X_1+X_2}(y) = \int_{-\infty}^{\infty} f_{X_1,\,X_2}(x_1,\, y-x_1)\, dx_1
            .\] Se inoltre $X_1,\,X_2$ sono \underline{indipendenti}, possiamo scrivere:
            \begin{equation}\label{eq:Somma_variabili_aleatorie_continue}
                f_{X_1+X_2}(y) = \int_{-\infty}^{\infty} f_{X_1}(x_1)\cdot f_{X_2}(y-x_1)\, dx_1
            .\end{equation}
        \end{prty}
        \begin{obsv}[Somma di gaussiane indipendenti]\label{obsv:Somma_gaussiane_indipendenti}
            Consideriamo due variabili aleatorie gaussiane indipendenti a media nulla:
            \begin{align*}
                Z_1 \sim \mathcal{N}(0,\, \sigma^2_1) & & Z_2 \sim \mathcal{N}(0,\, \sigma^2_2)
            ;\end{align*}
            Applicando l'equazione \eqref{eq:Somma_variabili_aleatorie_continue} otteniamo:
            \begin{align*}
                f_{Z_1+Z_2}(y) &= \int_{-\infty}^{\infty} \frac{1}{2\pi \sigma_1 \sigma_2} \cdot e^{-\frac{x^2}{2\sigma_1^2}-\frac{(y-x)^2}{2\sigma_2^2}}\, dx \\
                               &= \frac{1}{\sqrt{2\pi(\sigma_1^2 + \sigma_2^2)}} \cdot e^{-\frac{y^2}{2(\sigma_1^2 + \sigma_2^2)}} \cdot \int_{-\infty}^{\infty} \frac{1}{\sqrt{2\pi \tau}}\cdot e^{-\frac{(x-\nu)^2}{2\tau}}\, dx
            ,\end{align*}
            dove valgono le seguenti considerazioni:
            \begin{align*}
                \nu \coloneqq \frac{y \cdot \sigma_1^2}{\sigma_1^2 + \sigma_2^2}; & &
                \tau \coloneqq \frac{\sigma_1^ \cdot \sigma_2^2}{\sigma_1^2 + \sigma_2^2}; & &
                \int_{-\infty}^{\infty} \frac{1}{\sqrt{2\pi \tau}}\cdot e^{-\frac{(x-\nu)^2}{2\tau}}\, dx = 1
            .\end{align*}
            Allora possiamo riscrivere la densità come: \[
                f_{Z_1+Z_2}(y) = \frac{1}{\sqrt{2\pi (\sigma_1^2 + \sigma_2^2)}} \cdot e^{-\frac{y^2}{2(\sigma_1^2 + \sigma_2^2)}}
            ;\] possiamo indicare tale densità anche nel modo seguente: \[
            Z_1 + Z_2 \sim \mathcal{N}(0,\, \sigma_1^2 + \sigma_2^2)
            .\] Se ora prendiamo due variabili aleatorio $X_1,\,X_2$ a media non nulla $\mu_1,\,\mu_2$, la densità della loro somma è una trasformazione lineare della precedente: \[
            X_1 + X_2 \sim (Z_1 + Z_2) + (\mu_1 + \mu_2) \sim \mathcal{N}(\mu_1 + \mu_2,\, \sigma_1^2 + \sigma_2^2)
            .\] Applicando questa osservazione a $n$ variabili aleatorie con densità gaussiana otteniamo:
            \begin{equation*}
                \sum_{i=1}^{n} X_1 \sim \mathcal{N}\left(\sum_{i=1}^{n} \mu_i,\, \sum_{i=1}^{n} \sigma_i^2\right)
            .\end{equation*}
        \end{obsv}
        \begin{defn}[Densità Gamma]\label{defn:Densità_Gamma}
            Siano $X_1,\, X_2$ due variabili aleatorie indipendenti, entrambe esponenziali di parametro $\mu > 0$; allora applicano l'equazione \eqref{eq:Somma_variabili_aleatorie_continue} possiamo scrivere la densità della somma $X_1+X_2$ come segue: \[
                f_{X_1+X_2}(y) = \begin{cases}
                    \int_{0}^{y} \mu e^{-\mu u} \cdot \mu e^{-\mu(y-u)}\, du = \mu^2 e^{-\mu u} y & \text{se $y > 0$;} \\
                    0 & \text{se $y \leq 0$.}
                \end{cases}
            \] Procedendo per induzione su $n$ possiamo affermare che $X_1,\, \ldots,\, X_{n}$ sono variabili aleatorie indipendenti, ciascuna con densità esponenziale di parametro $\mu$; la densità della loro somma si scrive come:
            \begin{equation}\label{eq:Densità_Gamma}
                f_{X_1+\ldots+X_n}(x) = \frac{\mu^n}{(n-1)!}x^{n-1} \cdot 
                e^{-\mu x} \cdot \mathbf{1}_{(0,\,+\infty)}(x)
            .\end{equation}
            Chiamiamo \eqref{eq:Densità_Gamma} \emph{densità Gamma} di parametri $n$ e $\mu$, e la indichiamo con: \[
                X_1 + \ldots + X_{n} \sim \Gamma(n,\, \mu)
            .\]
        \end{defn}
        \begin{prty}[Densità di funzione continua]\label{prty:Densità_funzione_continua}
            Consideriamo $U,\,V$ due aperti di $\mathbb{R}^n$, 
            un'applicazione biunivoca $\vec{g}\,:\, U \mapsto V$ differenziabile con continuità, 
            e ammetta inversa $\vec{g}^{\,-1}$ anch'essa differenziabile con continuità; 
            prendiamo inoltre il vettore aleatorio $\vec{X}$ $n$ \nbdash dimensionale assolutamente continuo, 
            con densità $f_{.,}$ per cui valga $P(X \in U) = 1$. Allora $\vec{Y} \coloneqq \vec{g}(\vec{X})$ è un
            vettore aleatorio assolutamente continuo con densità:
            \begin{equation}\label{eq:Densità_funzione_continua}
                f_{\vec{Y}}(\vec{y}) = \mathbf{1}_{V}(\vec{y}) \cdot 
                f_{\vec{X}}\left(\vec{g}^{\,-1}(\vec{y})\right) \cdot 
                \left|\det\left(\vec{J}\left(\vec{g}^{\,-1}(\vec{y})\right)\right)\right|
            ,\end{equation}
            dove $\vec{J}(\vec{g}^{\,-1}(\vec{y}))$ indica la \emph{matrice Jacobiana} associata alla 
            funzione $\vec{g}^{\,-1}$, valutata in $\vec{y}$: \[
                \vec{J}(\vec{g}^{\,-1}) = \begin{bmatrix}
                    \frac{\partial g_1^{-1}}{\partial y_1} & \frac{\partial g_1^{-1}}{\partial y_2} & \cdots & \frac{\partial g_1^{-1}}{\partial y_n}\\
                    \frac{\partial g_2^{-1}}{\partial y_1} & \frac{\partial g_2^{-1}}{\partial y_2} & \cdots & \frac{\partial g_2^{-1}}{\partial y_n}\\
                    \vdots & \vdots & \cdots & \vdots\\
                    \frac{\partial g_n^{-1}}{\partial y_1} & \frac{\partial g_n^{-1}}{\partial y_2} & \cdots & \frac{\partial g_n^{-1}}{\partial y_n}
                \end{bmatrix}
            .\]
        \end{prty}
        \begin{defn}[Trasformazioni affini di vettori continui]
            Consideriamo la matrice $\vec{A}$ quadrata di ordine $n$ e invertibile, 
            il vettore colonna $\vec{b} \in \mathbb{R}^n$, il vettore aleatorio colonna assolutamente 
            continuo $\vec{X}$ a $n$ dimensioni e con densità $f_{\vec{X}}$, e valga 
            $\vec{Y} = \vec{A} \cdot \vec{X} + \vec{b}$.

            Calcoliamo la densità di $Y$ applicando \eqref{eq:Densità_funzione_continua} alla trasformazione 
            biunivoca $\vec{g}\,:\, \mathbb{R}^n \mapsto \mathbb{R}^n$ tale che $\vec{g}(\vec{x}) = \vec{A} \cdot \vec{x} + \vec{b}$ e ammetta inversa $\vec{g}^{\,-1}(\vec{y}) = \vec{A}^{\,-1}(\vec{y} - \vec{b})$:
            \begin{equation}\label{eq:Trasformazioni_vettori_continui}
                f_{\vec{Y}}(\vec{y}) = f_{\vec{X}}\left(\vec{A}^{\,-1}(\vec{y} - \vec{b})\right) \cdot 
                \left|\det(\vec{A}^{\,-1})\right|
            .\end{equation}
        \end{defn}
    \section{Valore atteso per funzioni di vettori aleatori}
        \begin{defn}[discreto]
            Consideriamo un vettore aleatorio discreto $\vec{X}$ che assume valori in $S$ e 
            ha densità $p_{\vec{X}}$, una funzione $g\,:\, \mathbb{R}^n \mapsto \mathbb{R}$ e una 
            variabile aleatoria $Y \coloneqq g(\vec{X})$.

            Se si verifica: \[
                \sum_{\vec{x} \in S} \left|g(\vec{x})\right| \cdot p_{\vec{X}}(\vec{x}) < +\infty
            ,\] allora $Y$ ammette valore atteso ed esso si calcola come:
            \begin{equation}\label{eq:Valore_atteso_vettori_discreti}
                \text{E}(Y) = \sum_{\vec{x} \in S} g(\vec{x}) \cdot p_{\vec{X}}(\vec{x})
            .\end{equation}
        \end{defn}
        \begin{defn}[continuo]
            Consideriamo un vettore aleatorio assolutamente continuo $\vec{X}$ che abbia densità 
            $f_{\vec{X}}$, una funzione $g\,:\, \mathbb{R}^n \mapsto \mathbb{R}$ e una variabile aleatoria 
            $Y \coloneqq g(\vec{X})$.

            Se si verifica: \[
                \int_{\mathbb{R}^n} \left|gx_1,\, \ldots,\, x_{n})\right| \cdot f_{\vec{X}}(x_1,\, \ldots,\, x_{n})\, dx_1 \cdot \ldots \cdot dx_n < +\infty
            ,\] allora $Y$ ammette valore atteso ed esso si calcola come:
            \begin{equation}\label{eq:Valore_atteso_vettori_continui}
                \text{E}(Y) = \int_{\mathbb{R}^n} g(x_1,\, \ldots,\, x_{n}) \cdot f_{\vec{X}}(x_1,\, \ldots,\, x_{n})\, dx_1 \cdot \ldots \cdot dx_n
            .\end{equation}
        \end{defn}
    \section{Proprietà del valore atteso per due vettori aleatori}
    \begin{prty}[Media della somma]
            Siano $X_1,\,X_2$ variabili aleatorie definite sullo stesso spazio di probabilità, ed entrambe ammettano media; allora anche $X_1 + X_2$ ammette media ed essa vale: \[
            \text{E}(X_1 + X_2) = \text{E}(X_1) + \text{E}(X_2)
            .\]
        \end{prty}
        \begin{proof}
            Prendiamo il vettore aleatorio assolutamente continuo $(X_1,\,X_2)$ con densità $f_{X_1,\,X_2}$; 
            applicando la disuguaglianza triangolare $|x+y| \leq |x| + |y|$ otteniamo:
            \begin{align*}
                \int_{\mathbb{R}^2} |x+y| \cdot f_{X_1,\,X_2}(x,\,y)\, dx\,dy &\leq 
                \int_{\mathbb{R}^2} |x| \cdot f_{X_1,\,X_2}(x,\,y)\, dx\,dy + 
                \int_{\mathbb{R}^2} |y| \cdot f_{X_1,\,X_2}(x,\,y)\, dx\,dy \\
                &= \int_{\mathbb{R}} |x|\cdot f_{X_1}(x)\, dx + \int_{\mathbb{R}} |y|\cdot f_{X_2}(y)\, dy \\
                &< \infty
            ;\end{align*}
            quindi $X_1 + X_2$ ammette media, e applicando \eqref{eq:Valore_atteso_vettori_continui} alla 
            funzione $g(x,\,y) = x+y$ otteniamo:
            \begin{align*}
                \text{E}(X_1 + X_2) &= \int_{\mathbb{R}^2} (x+y)\cdot f_{X_1,\,X_2}(x,\,y)\, dx\, dy \\
                &= \int_{\mathbb{R}} x\cdot \left(\int_{\mathbb{R}}f_{X_1,\,X_2}(x,\,y)\, dy\right)\, dx + \int_{\mathbb{R}} y\cdot \left(\int_{\mathbb{R}}f_{X_1,\,X_2}(x,\,y)\, dx\right)\, dy \\
                &= \int_{\mathbb{R}} x\cdot f_{X_1}(x)\, dx + \int_{\mathbb{R}} y\cdot f_{X_2}(y)\, dy \\
                &= \text{E}(X_1) + \text{E}(X_2) \qedhere
            .\end{align*}
        \end{proof}
        \begin{obsv}
            Possiamo estendere la proprietà della somma dei valori attesi ad $n$ variabili aleatorie: \[
            \text{E}(X_1+\ldots+X_n) = \text{E}(X_1)+\ldots+\text{E}(X_n)
            .\] 
        \end{obsv}
        \begin{prty}[Media del prodotto]
            Siano $X_1,\,X_2$ variabili aleatorie indipendenti, ed entrambe ammettano media; allora anche $X_1 \cdot X_2$ ammette media ed essa vale:
            \begin{equation}\label{eq:Media_prodotto_v_a}
                \text{E}(X_1 \cdot X_2) = \text{E}(X_1) \cdot \text{E}(X_2)
            .\end{equation}
        \end{prty}
        \begin{proof}
            Prendiamo due variabili aleatorie continue $X_1,\,X_2$ con densità $f_{X_1},\, f_{X_2}$; possiamo scrivere la densità congiunta rispetto alle marginali: \[
                \iint_{\mathbb{R}\times\mathbb{R}} |x\cdot y| \cdot f_{X_1}(x) \cdot f_{X_2}(y)\, dx\, dy =
                \int_{\mathbb{R}} |x|\cdot f_{X_1}(x)\, dx \cdot \int_{\mathbb{R}} |y|\cdot f_{X_2}(y)\, dy 
                < +\infty
            .\] Possiamo affermare che esiste $\text{E}(X_1 \cdot X_2)$ per \eqref{eq:Valore_atteso_vettori_continui} applicata alla funzione $g(x,\,y) = x\cdot y$; inoltre possiamo scrivere:
            \begin{align*}
                \text{E}(X_1 \cdot X_2) &=
                \iint_{\mathbb{R}\times\mathbb{R}} |x\cdot y| \cdot f_{X_1}(x) \cdot f_{X_2}(y)\, dx\, dy \\
                &= \int_{\mathbb{R}} x\cdot f_{X_1}(x)\, dx \cdot \int_{\mathbb{R}} y\cdot f_{X_2}(y)\, dy \\
                &= \text{E}(X_1) \cdot \text{E}(X_2) \qedhere
            .\end{align*}
        \end{proof}
        \begin{prty}[Varianza della somma]\label{prty:Varianza_somma_v_a}
            Siano $X_1,\,X_2$ variabili aleatorie che ammettano varianza finita $\text{Var}(X_1),\,\text{Var}(X_2)$; allora anche $X_1 + X_2$ ammette varianza finita, ed essa vale: \[
            \text{Var}(X_1 + X_2) = \text{Var}(X_1) + \text{Var}(X_2) + 
            2 \text{E}\left[(X_1 - \text{E}(X_1))\cdot(X_2 - \text{E}(X_2))\right]
            .\] Se inoltre $X_1,\,X_2$ sono \underline{indipendenti}, allora otteniamo:
            \begin{equation}\label{eq:Varianza_somma_v_a}
                \text{Var}(X_1 + X_2) = \text{Var}(X_1) + \text{Var}(X_2)
            .\end{equation}
        \end{prty}
        \begin{proof}
            Sappiamo che valgono le seguenti reazioni col valore atteso della somma di due variabili aleatorie:
            \begin{align*}
                \left((X_1 + X_2) - \text{E}(X_1 + X_2)\right)^2 &= \left[(X_1 - \text{E}(X_1)) + (X_2 + \text{E}(X_2))\right]^2 \\
                                      &\leq 2\left[(X_1 - \text{E}(X_1))^2 + (X_2 + \text{E}(X_2))^2\right]
            ;\end{align*}
            allora applicando le definizione di varianza possiamo scrivere:
            \begin{align*}
                \text{Var}(X_1 + X_2) &= \text{E}\left[((X_1+X_2)-\text{E}(X_1+X_2))^2\right] \\
                &= \text{E}\left[((X_1 - \text{E}(X_1))+(X_2 - \text{E}(X_2)))^2\right] \\
                &\leq 2\left[\text{E}(X_1 - \text{E}(X_1))^2+\text{E}(X_2-\text{E}(X_2))^2\right] \\
                &= 2\left[\text{Var}(X_1)+\text{Var}(X_2)\right]
            .\end{align*}
            Quindi se $X_1,\,X_2$ ammettono varianza, anche $X_1+X_2$ la ammette; essa si scrive nel modo 
            seguente:
            \begin{align*}
                \text{Var}(X_1+X_2) &= \text{E}\left[(X_1 - \text{E}(X_1))^2\right] + 
                \text{E}\left[(X_2 - \text{E}(X_2))^2\right] + 
                2\, \text{E}\left[(X_1 - \text{E}(X))\cdot(X_2 - \text{E}(X_2))\right] \\
                &= \text{Var}(X_1) + \text{Var}(X_2) + 
                2\, \text{E}\left[(X_1 - \text{E}(X))\cdot(X_2 - \text{E}(X_2))\right]
            .\end{align*}
            Infine notiamo che , se $X_1,\,X_2$ sono indipendenti, allora per il terzo termine vale: \[
                2\, \text{E}\left[(X_1 - \text{E}(X))\cdot(X_2 - \text{E}(X_2))\right] = 0
            ,\] e otteniamo proprio l'equazione \eqref{eq:Varianza_somma_v_a}: \[
                \text{Var}(X_1 + X_2) = \text{Var}(X_1) + \text{Var}(X_2). \qedhere
            \]
        \end{proof}
        \begin{obsv}
            Possiamo estendere la proprietà della somma delle varianze ad $n$ variabili aleatorie:
            \begin{equation}\label{eq:Somma_varianza_v_a}
                \text{Var}(X_1+\ldots+X_n) = \sum_{i=1}^{n} \text{Var}(X_i) + 
                2 \sum_{i=1}^{n-1} \sum_{j=i+1}^{n} \text{E}\left[(X_i-\text{E}(X_i))\cdot(X_j-\text{E}(X_j))\right]
            .\end{equation}
        \end{obsv}
    \section{Media e varianza della Binomiale}
        \begin{obsv}
            Consideriamo $X\sim \mathcal{B}i(n,\,p)$; essa avrà la stessa densità della somma di $n$ variabili aleatorie indipendenti ($X_1,\, \ldots,\, X_{n}$) con densità di Bernoulli di parametro $p$; allora ricaviamo:
            \begin{align*}
                \text{E}(X) = \text{E}\left(\sum_{i=1}^{n} X_i\right) = \sum_{i=1}^{n} \text{E}(X_i) =
                \sum_{i=1}^{n} p = n \cdot p
                & & \text{per Proprietà}~\ref{eq:Valore_atteso_vettori_continui}
            ;\end{align*}
            \begin{align*}
                \text{Var}(X) = \text{Var}(\sum_{i=1}^{n} X_i) = \sum_{i=1}^{n} \text{Var}(X_i) = 
                \sum_{i=1}^{n} p(1-p) = n \cdot p(1-p)
                & & \text{per equazione}~\eqref{eq:Somma_varianza_v_a}
            .\end{align*}
        \end{obsv}
    \section{Covarianza}
        \begin{defn}
            Siano $X_1,\,X_2$ variabili aleatorie definite sullo stesso spazio di probabilità ed entrambe ammettano varianza; allora chiamiamo \emph{covarianza} di $X_1$ e $X_2$ il numero calcolato come:
            \begin{equation}\label{eq:Covarianza}
                \text{Cov}(X_1,\,X_2) \coloneqq \text{E}\left[(X_1-\text{E}(X_1))\cdot(X_2-\text{E}(X_2))\right]
            .\end{equation}
        \end{defn}
        \begin{obsv}
            Estendendo la Proprietà~\ref{prty:Varianza_somma_v_a} con la definizione di covarianza \eqref{eq:Covarianza}, otteniamo: \[
            \text{Var}(X_1+X_2) = \text{Var}(X_1) + \text{Var}(X_2) + 2\,\text{Cov}(X_1,\,X_2)
        ;\] inoltre intuiamo da \eqref{eq:Varianza_somma_v_a} che per variabili aleatorie indipendenti la covarianza è nulla.
        \end{obsv}
        \begin{prty}
            Siano $X_1,\,X_2,\,X_3$ variabili aleatorie con varianza finita e $a,\,b \in \mathbb{R}$; possiamo affermare che:
            \begin{enumerate}
                \item $\text{Cov}(X_1,\,X_2) = \text{Cov}(X_2,\,X_1)$;
                \item $\text{Cov}(X_1,\,X_2) = \text{E}(X_1\cdot X_2) - \text{E}(X_1)\cdot \text{E}(X_2)$;
                \item $\text{Cov}(a\cdot X_1,\,X_2) = a\cdot \text{Cov}(X_1,\,X_2)$;
                \item $\text{Cov}(X_1+X_2,\,X_3) = \text{Cov}(X_1,\,X_3)+\text{Cov}(X_2,\,X_3)$;
                \item se $X_1,\,X_2$ sono indipendenti allora $\text{Cov}(X_1,\,X_2)=0$ (non vale l'inverso).
            \end{enumerate}
        \end{prty}
        \begin{proof}
            \hfill
            \begin{enumerate}
                \item La definizione di Covarianza contiene il valore atteso del prodotto di $(X_1-\text{E}(X_1))$ con $(X_2-\text{E}(X_2))$; quest'ultimo gode della proprietà commutativa, per cui:
                    \begin{align*}
                        \text{Cov}(X_1,\,X_2) &\coloneqq 
                        \text{E}\left[(X_1-\text{E}(X_1))\cdot(X_2-\text{E}(X_2))\right] \\
                        &\,= \text{E}\left[(X_2-\text{E}(X_2))\cdot(X_1-\text{E}(X_1))\right] =
                        \text{Cov}(X_2,\,X_1)
                    .\end{align*}
                \item Dalla definizione di covarianza espandiamo il prodotto al esondo membro:
                    \begin{align*}
                        \text{Cov}(X_1,\,X_2) &= \text{E}(X_1\cdot X_2 - \text{E}(X_1)\cdot X_2 -\text{E}(X_2)\cdot X_1 +\text{E}(X_2)\cdot \text{E}(X_1)) \\
                        &= \text{E}(X_1\cdot X_2) - \text{E}(X_1)\cdot \text{E}(X_2) \cancel{-\text{E}(X_1)\cdot \text{E}(X_2)} \cancel{+ \text{E}(X_1)\cdot \text{E}(X_2)} \\
                        &= \text{E}(X_1 \cdot X_2) - \text{E}(X_1)\cdot \text{E}(X_2)
                    .\end{align*}
                \item Usando la formula appena dimostrata, per il punto \eqref{itm:Linearità_valore_atteso_somma} della Proprietà~\ref{prty:Valore_atteso} otteniamo:
                    \begin{align*}
                        \text{Cov}(a\cdot X_1,\,X_2) &= \text{E}(a\cdot X_1 \cdot X_2) - \text{E}(a\cdot X_1)\cdot \text{E}(X_2) \\
                        &= a\cdot\text{E}(X_1 \cdot X_2) - a\cdot \text{E}(X_1) \cdot \text{E}(X_2) \\
                        &= a \cdot \left[\text{E}(X_1 \cdot X_2) - \text{E}(X_1) \cdot \text{E}(X_2)\right] \\
                        &= a \cdot \text{Cov}(X_1,\,X_2)
                    .\end{align*}
                \item Sviluppando il prodotto con $(X_1 + X_2)$ otteniamo:
                    \begin{align*}
                        \text{Cov}(X_1+X_2,\,X_3) &= \text{E}((X_1+X_2)\cdot X_3) - \text{E}(X_1+X_2) \cdot \text{E}(X_3) \\
                        &= \text{E}(X_1 \cdot X_3 + X_2 \cdot X_3) - \left[\text{E}(X_1)+\text{E}(X_2)\right]\cdot \text{E}(X_3) \\
                        &= \text{E}(X_1 \cdot X_3) - \text{E}(X_1)\cdot \text{E}(X_2) + \text{E}(X_2 \cdot X_3) - \text{E}(X_2)\cdot \text{E}(X_3) \\
                        &= \text{Cov}(X_1,\,X_3) + \text{Cov}(X_2,\,X_3) \qedhere
                    .\end{align*}
            \end{enumerate}
        \end{proof}
    \section{Coefficiente di correlazione lineare}
        \begin{defn}
            Prendiamo due variabili aleatorie $X_1,\,X_2$ definite sullo stesso spazio di probabilità, che ammettano entrambe covarianza ($0<\text{Var}(X_1) \land 0<\text{Var}(X_2)$); allora definiamo il \emph{coefficiente di correlazione} di $X_1,\,X_2$ il numero calcolato come: \[
            \rho_{X_1,\,X_2} \coloneqq \frac{\text{Cov}(X_1,\,X_2)}{\sqrt{\text{Var}(X_1)\cdot \text{Var}(X_2)}}
            .\] 
        \end{defn}
        \begin{prty}
            Siano $X_1,\,X_2$ variabili aleatorie con varianza finita e $a..b \in \mathbb{R}$; allora vale $|\rho_{X_1,\,X_2}| \leq 1$ e inoltre: \[
                |\rho_{X_1,\,X_2}| = 1 \iff \exists\, a,\,b \in \mathbb{R}\,:\, P(X_2=a\cdot X_1+b) = 1
            ;\] se ciò si verifica, allora i due reali valgono:
            \begin{align*}
                a = \frac{\text{Cov}(X_1,\,X_2)}{\text{Var}(X_1)}; & &
                b = \text{E}(X_2) - \frac{\text{E}(X_1)\cdot \text{Cov}(X_1,\,X_2)}{\text{Var}(X_1)}
            .\end{align*}
        \end{prty}
        \begin{proof}
            Chiamiamo $\text{Var}(X_1) = \sigma_1^2$ e $\text{Var}(X_2) = \sigma_2^2$; usando le proprietà della varianza scriviamo:
            \begin{align*}
                0 &\leq \text{Var}\left(\frac{X_1}{\sigma_1}+\frac{X_2}{\sigma_2}\right)=
                \frac{\text{Var}(X_1)}{\sigma_1^2} + \frac{\text{Var}(X_2)}{\sigma_2^2} + 2\,\text{Cov}\left(\frac{X_1}{\sigma_1},\,\frac{X_2}{\sigma_2}\right) \\
                &= \frac{\sigma_1^2}{\sigma_1^2} + \frac{\sigma_2^2}{\sigma_2^2} + 2\,\frac{\text{Cov}(X_1,\,X_2)}{\sigma_1 \cdot \sigma_2} \\
                &= 2\, (1 + \rho_{X_1,\,X_2}) \implies \rho_{X_1,\,X_2} \geq -1
            .\end{align*}
            Vale inoltre:
            \begin{align*}
                0 &\leq \text{Var}\left(\frac{X_1}{\sigma_1}-\frac{X_2}{\sigma_2}\right) = 
                \frac{\text{Var}(X_1)}{\sigma_1^2} + \frac{\text{Var}(X_2)}{\sigma_2^2} 
                -2\, \frac{\text{Cov}(X_1,\,X_2)}{\sigma_1 \cdot \sigma_2} \\
                  &= 2\, (1 - \rho_{X_1,\,X_2}) \implies \rho_{X_1,\,X_2} \leq 1
            .\end{align*}
            Dal momento che: \[
                \rho_{X_1,\,X_2} = 1 \iff \text{Var}\left(\frac{X_1}{\sigma_1} - \frac{X_2}{\sigma_2}\right) = 0
            ,\] segue dalle proprietà della varianza: \[
            \rho_{X_1,\,X_2} = 1 \iff P\left(\frac{X_1}{\sigma_1}-\frac{X_2}{\sigma_2} = 
            \frac{\text{E}(X_1)}{\sigma_1}-\frac{\text{E}(X_2)}{\sigma_2}\right) = 1
            ;\] inoltre sappiamo che: \[
            \rho_{X_1,\,X_2} = 1 \iff \text{Cov}(X_1,\,X_2)=\sigma_1 \cdot \sigma_2
            ,\] da cui segue: \[
            X_2 = \text{E}(X_2) + \frac{\text{Cov}(X_1,\,X)}{\sigma_1^2} \cdot (X_1 - \text{E}(X_1))
            .\] Invece per $\rho_{X_1,\,X_2} = -1$ otteniamo:
            \begin{align*}
                \rho_{X_1,\,X_2} = -1 &\iff \text{Cov}(X_1,\,X_2) = -\sigma_1 \cdot \sigma_2 \\
                &\iff \text{Var}\left(\frac{X_1}{\sigma_1}+\frac{X_2}{\sigma_2}\right) = 0 \\
                &\iff P\left(\frac{X_1}{\sigma_1}+\frac{X_2}{\sigma_2}= \frac{\text{E}(X_1)}{\sigma_1}+\frac{\text{E}(X_2)}{\sigma_2}\right) = 1 \\
                &\iff X_2 = \text{E}(X_2) + \frac{\text{Cov}(X_1,\,X_2)}{\sigma_1^2} \cdot (X_1 - \text{E}(X_1))
                \qedhere
            .\end{align*}
        \end{proof}
        \begin{note}
            Esiste sempre un legame lineare tra due variabili aleatorie $X_1,\,X_2$ definito come $X_2 = a\cdot X_1 + b$, se e solo se $\rho_{X_1,\,X_2}=\pm 1$; inoltre valgono le seguenti implicazioni:
            \begin{itemize}
                \item $\rho_{X_1,\,X_2} = -1 \implies a<0 \land \text{Cov}(X_1,\,X_2)<0$;
                \item $\rho_{X_1,\,X_2} = 1 \implies a>0 \land \text{Cov}(X_1,\,X_2)>0$.
            \end{itemize}
        \end{note}
    %\section{Distribuzione del massimo e del minimo}
        %TODO: Completare sezione
    \section{Media campionaria}\label{sec:Media_campionaria}
        \begin{defn}
            Consideriamo una popolazione di $n$ elementi, a ciascuno dei quali è associata una grandezza numerica $X_i$ con media $\mu$ e varianza $\sigma^2$; l'insieme di queste grandezze (variabili aleatorie indipendenti identicamente distribuite) $X_1,\, \ldots,\, X_{n}$ è detto \emph{campione}; chiamiamo \emph{media campionaria} la sua media, e la indichiamo come: \[
            \overline{X} \coloneqq \frac{X_1 + \ldots + X_{n}}{n}
            .\] Dato che $\overline{X}$ è una funzione di variabili aleatorie indipendenti identicamente distribuite, anch'essa è una variabile aleatoria, quindi ammette media e varianza nel modo seguente:
            \begin{itemize}
                \item Il valore atteso della media campionaria vale:
                    \begin{align*}
                        \text{E}(\overline{X}) &= \text{E}\left(\frac{X_1 + \ldots + X_n}{n}\right)
                        = \frac{\text{E}(X_1) + \ldots + \text{E}(X_n)}{n} \\
                        &= \frac{n \cdot \mu}{n} = \mu
                    .\end{align*}
                \item La varianza della media campionaria vale:
                    \begin{align*}
                        \text{Var}(\overline{X}) &= \text{Var}\left(\frac{X_1 + \ldots + X_n}{n}\right)
                        = \frac{\text{Var}(X_1) + \ldots + \text{Var}(X_n)}{n} \\
                        &= \frac{n\cdot \sigma^2}{n^2} = \frac{\sigma^2}{n}
                    .\end{align*}
            \end{itemize}
        \end{defn}
    \section{Legge dei grandi numeri}\label{sec:Legge_grandi_numeri}
        \begin{thm}[Legge debole]\label{thm:Legge_grandi_numeri_debole}
            Sia $X_1,\, \ldots,\, X_{n}$ una successione di variabili aleatorie indipendenti identicamente distribuite con media $\mu$ e varianza $\sigma^2$ finite, inoltre $S_n = X_1 + \ldots + X_n$; allora vale: \[
                \forall \varepsilon > 0 \,:\, \lim_{n \to \infty} P\left(\left|\frac{S_n}{n}-\mu\right| > \varepsilon\right) = 0
            .\]
        \end{thm}
        \begin{proof}
            Dato che le variabili $X_i$ sono indipendenti identicamente distribuite otteniamo:
            \begin{align*}
                \text{Var}(S_n) = n\cdot \text{Var}(X_1) = n\cdot \sigma^2  \implies
                \text{Var}\left(\frac{S_n}{n}\right) = \frac{n\, \sigma^2}{n^2} = \frac{\sigma^2}{n};
                & & \text{E}\left(\frac{S_n}{n}\right) = \mu
            .\end{align*}
            Dalla disuguaglianza di Chebychev \eqref{eq:Disuguaglianza_Chebychev} segue che: \[
                \forall \varepsilon > 0 \,:\, \lim_{n \to \infty} \left[P\left(\left|\frac{S_n}{n}-\mu\right|> \varepsilon\right) \leq \frac{\sigma^2}{n\, \varepsilon^2}\right] \rightarrow 0. \qedhere
            \]
        \end{proof}
        \begin{thm}[Legge forte]\label{thm:Legge_grandi_numeri_forte}
            Sia $X_1,\, \ldots,\, X_{n}$ una successione di variabili aleatorie indipendenti identicamente distribuite con media finita $\mu$; allora possiamo affermare che: \[
                P\left(\left\{\omega\,:\, \lim_{n \to \infty} \frac{S_n(\omega)}{n} = \mu\right\}\right) = 1
            .\]
        \end{thm}
        \begin{prty}
            Sia $h$ una funzione su $[0,\,1]$ per cui valga $\int_{[0,\,1]} |h(x)|\, dx < +\infty$, 
            prendiamo inoltre $U_1,\, U_2,\, \ldots$ variabili aleatorie indipendenti 
            identicamente distribuite con densità uniforme su $[0,\,1]$; allora: \[
                P\left(I_{\mathbf{1}n} \coloneqq \frac{1}{n} \cdot \sum_{j=1}^{n} h(U_j) \rightarrow
                \int_{0}^{1} h(x)\, dx \text{ per } n \rightarrow \infty\right) = 1
            .\] 
        \end{prty}
        \begin{proof}
            Dato che le variabili aleatorie $h(U_1),\, h(U_2),\, \ldots$ sono indipendenti identicamente distribuite con media finita, applicando la \emph{legge dei grandi numeri}.
        \end{proof}
    \section{Teorema del limite centrale} 
        \begin{obsv}
            Consideriamo $n$ variabili aleatorie $X_1,\, \ldots,\, X_{n}$ indipendenti identicamente distribuite con media $\mu$ e varianza $\sigma^2$ \underline{finite}; la legge dei grandi numeri ci permette di affermare (per $n$ grande) \[
            \overline{X}_n \simeq \mu
            .\] Se inoltre le variabili $X_i$ sono gaussiane, si ha: \[
            \overline{X}_n \sim \mathcal{N}\left(\mu,\,\frac{\sigma^2}{n}\right) \implies
                \frac{\sqrt{n}}{\sigma} \cdot (\overline{X}_n - \mu) \sim \mathcal{N}(0,\,1) 
            ,\] da cui otteniamo: \[
            P\left(|\overline{X}_n - \mu| \leq \delta\right) = 
            \Phi\left(\frac{\sqrt{n}}{\sigma}\,\delta\right) - 
            \Phi\left(-\frac{\sqrt{n}}{\sigma}\,\delta\right) =
            2\, \Phi\left(\frac{\sqrt{n}}{\sigma}\,\delta\right) - 1
            .\]
        \end{obsv}
        \begin{thm}[Limite centrale]\label{thm:Teorema_limite_centrale}
            Sia $X_1,\, X_2,\, \ldots$ una successione di variabili aleatorie indipendenti identicamente distribuite con media $\mu$ e varianza $\sigma^2$, per cui valga $0 < \sigma^2 < +\infty$; allora possiamo scrivere:
            \begin{equation}\label{eq:Teorema_limite_centrale}
                \forall x\ in \mathbb{R} \,:\, \lim_{n \to \infty} P\left(\frac{\sqrt{n}}{\sigma} \cdot (\overline{X}_n - \mu) \leq x\right) = \int_{-\infty}^{x} \frac{1}{\sqrt{2\pi}}\cdot e^{-u^2 /2}\, du = \Phi(x)
            .\end{equation}
        \end{thm}
        \begin{obsv}
            Prendendo un $n$ adeguatamente elevato, la ripartizione della standardizzata della media campionaria è approssimabile con con la ripartizione gaussiana standard.
        \end{obsv}

    %! TEX root = main.tex
% Capitolo 4

\chapter{Distribuzioni notevoli}
    \section{Matrice delle covarianza}\label{defn:Matrice_covarianza}
        \begin{defn}
            Consideriamo un vettore aleatorio $n$\nbdash dimensionale $\vec{X} = (X_1,\, \ldots,\, X_{n})$ tale che sia definita la varianza di ciascuna delle sue componenti; allora chiamiamo \emph{matrice di covarianza} di $\vec{X}$ la matrice quadrata di ordine $n$ definita come:
            \begin{align*}
                \forall i,\,j \in [1,\,n] \,:\, \vec{C}_{\vec{X}} \coloneqq [c_{i,\,j}] & &
                \text{per } c_{i,\,j} = \text{Cov}(X_i,\, X_j)
            .\end{align*}
        \end{defn}
        \begin{prty}
            Sia $\vec{C}_{\vec{X}}$ la matrice di covarianza di un vettore aleatorio $\vec{X}$; allora possiamo affermare che:
            \begin{enumerate}
                \item $\vec{C}_{\vec{X}}$ è una matrice simmetrica e semi-definita positiva 
                    ($\forall \vec{X} \neq 0 \in \mathbb{R}^n \,:\, 
                    \vec{X}^{\text{\,T}}\cdot \vec{C}_{\vec{X}} \cdot \vec{X} \geq 0$);
                \item se abbiamo $\vec{A}$ una matrice $m \times n$ e $\vec{b}$ un vettore di dimensione $m$ allora la matrice di covarianza di $\vec{Y} = \vec{A} \cdot \vec{X} + \vec{b}$ si scrive come: \[
                \vec{C}_{\vec{Y}} = \vec{A} \cdot \vec{C}_{\vec{X}} \cdot \vec{A}^{\text{\,T}}
                .\] 
            \end{enumerate}
        \end{prty}
        \begin{proof}
            \hfill
            \begin{enumerate}
                \item La matrice di covarianza risulta simmetrica poiché, preso il suo elemento di posizione $i,\,j$ vale: \[
                c_{i,\,j} = \text{Cov}(X_{i},\,X_{j}) = \text{Cov}(X_{j},\,X_{i}) = c_{j,\,i}
            .\] Inoltre, se prendiamo un vettore colonna $\vec{\lambda} \in \mathbb{R}^n$ per $i\in [1,\,m]$ e $j\in [1,\,n]$ otteniamo:
                \begin{align*}
                    \vec{\lambda}^{\text{\,T}} \cdot \vec{C}_{\vec{X}} \cdot \vec{\lambda} &=
                    \sum_{i=1}^{n} \sum_{j=1}^{n}\, [\lambda_i \lambda_j \cdot c_{i,\,j}] \\
                    &= \text{E}\left(\sum_{i=1}^{n} \sum_{j=1}^{n} \lambda_i \lambda_j \cdot 
                    \left(X_i-\text{E}(X_i)\right)\left(X_j-\text{E}(X_j)\right)\right) \\
                    &= \text{E}\left(\sum_{i=1}^{n} \lambda_i \cdot \left(X_i-\text{E}(X_i)\right)^2\right) \geq 0
                .\end{align*}
            \item Se le componenti di $\vec{X}$ hanno varianza finita allora anche le componenti di $\vec{Y} = \vec{A} \cdot \vec{X} + \vec{b}$ la avranno; allora possiamo studiare la matrice di covarianza $\vec{C}_{\vec{Y}}$ considerando $\tilde{c}_{i,\,j}$ l'elemento di posto $(i,\,j)$ di $\vec{C}_{\vec{Y}}$, ottenendo:
                \begin{align*}
                    \tilde{c}_{i,\,j} = \text{Cov}\left(\vec{Y}_i,\, \vec{Y}_j\right)
                    &= \text{Cov}\left(\sum_{k=1}^{n} [a_{i,\,k} \cdot X_k + b_i],\, 
                    \sum_{\ell=1}^{n} [a_{j,\,\ell} \cdot X_\ell + b_j]\right) \\
                    &= \sum_{k=1}^{n} \sum_{\ell=1}^{n}\, \left[a_{i,\,k} \cdot a_{j,\,\ell} \cdot \text{Cov}(X_k,\, X_{\ell})\right]
                ,\end{align*}
                esso è proprio l'elemento $(i,\,j)$ della matrice $\vec{A} \cdot \vec{C}_{\vec{X}}\vec{A}^{\text{\,T}}$. \qedhere
            \end{enumerate}
        \end{proof}
    \section{Vettori gaussiani}
        \begin{defn}[Normale standard]\label{defn:Vettore_gaussiano_standard}
            Il vettore aleatorio $\vec{Z} = (Z_1,\, \ldots,\, Z_{n})^{\text{T}}$ si dice \emph{gaussiano standard} $n$\nbdash dimensionale se le variabili aleatorie $Z_1,\, \ldots,\, Z_{n}$ sono gaussiane standard indipendenti.

            Un vettore aleatorio $\vec{Z}$ gaussiano appena definito è \underline{assolutamente continuo} e ha densità: \[
                f_{\vec{Z}}(z_1,\, \ldots,\, z_{n}) = \frac{1}{(2\pi)^{n /2}} \cdot e^{-\frac{1}{2}\sum_{k=1}^{n} z_k^2}
            ;\] questo risultato si ottiene dall'indipendenza delle componenti di $\vec{Z}$: \[
            f_{\vec{Z}}(z_1,\, \ldots,\, z_{n}) = f_{Z_1}(z_1) \cdot \ldots \cdot f_{Z_n}(z_n) =
            \frac{1}{\sqrt{2\pi}} \cdot e^{- z_1^2 /2} \cdot \ldots \cdot \frac{1}{\sqrt{2\pi}} \cdot e^{- z_n^2 /2} = \frac{1}{(2\pi)^{n /2}} \cdot e^{-\frac{1}{2}\sum_{k=1}^{n} z_k^2}
            .\] 
        \end{defn}
        \begin{obsv}
            Media e matrice di covarianza di un vettore gaussiano standard si ottengono come:
            \begin{gather*}
                \text{E}(\vec{Z}) = \text{\O} = (0,\,0,\,\ldots); \\
                \vec{C}_{\vec{Z}} = \vec{I} \in \mathbb{R}^n \times \mathbb{R}^n.
            \end{gather*}
        \end{obsv}
        \begin{defn}[Normale]\label{defn:Vettore_gaussiano}
            Un vettore aleatorio $n$\nbdash dimensionale $\vec{X}$ è detto \emph{gaussiano} o normale se esistono una matrice $\vec{A} \in \mathbb{R}^n \times \mathbb{R}^m$, un vettore $\vec{\mu} \in \mathbb{R}^n$ e un vettore gaussiano \underline{standard} $\vec{Z} \in \mathbb{R}^m$ tali che: \[
            \vec{X} = \vec{A} \cdot \vec{Z} + \vec{\mu}
            .\] 
        \end{defn}
        \begin{obsv}\label{obsv:Vettore_gaussiano}
            Media e matrice di covarianza di un vettore gaussiano si ottengono come:
            \begin{gather*}
                \text{E}(\vec{X}) = \vec{A} \cdot \text{E}(\vec{Z}) + \vec{\mu} = \vec{\mu}; \\
                \vec{C}_{\vec{X}} = \vec{A} \cdot \vec{I} \cdot \vec{A}^{\text{\,T}} = \vec{A} \cdot \vec{A}^{\text{\,T}}.
            \end{gather*}
        \end{obsv}
        \begin{prty}
            Un vettore gaussiano $\vec{X} = \vec{A} \cdot \vec{Z} + \vec{\mu}$ ha densità su $\mathbb{R}^n$ se e solo se la matrice di covarianza $\vec{C} = \vec{A} \cdot \vec{A}^{\text{\,T}}$ è non-singolare; in tal caso la densità vale:
            \begin{equation}\label{eq:Densità_vettore_gaussiano}
                f_{\vec{X}}(\vec{x}) = \frac{1}{\sqrt{(2\pi)^n \cdot \det[\vec{C}]}} \cdot e^{-\frac{1}{2}(\vec{x}-\vec{\mu})^{\text{T}} \cdot \vec{C}^{-1}\cdot (\vec{x}-\vec{\mu})}
            .\end{equation}
        \end{prty}
        \begin{obsv}
            Consideriamo un vettore aleatorio assolutamente continuo $\vec{X}$ con densità \eqref{eq:Densità_vettore_gaussiano}, dove $\vec{C}$ è una matrice simmetrica definita positiva; allora possiamo estrarre la radice di $\vec{C}$ come: \[
            \exists \vec{A} \,:\, \vec{C} = \vec{A} \cdot \vec{A}^{\text{\,T}}
        .\] Prendiamo ora un vettore $\vec{Z} = \vec{A}^{\,-1}\cdot(\vec{X} - \vec{\mu})$, e calcoliamo la sua densità usando \eqref{eq:Trasformazioni_vettori_continui}:
        \begin{align*}
            f_{\vec{Z}}(\vec{z}) = f_{\vec{X}}(\vec{A} \cdot \vec{z} + \vec{\mu}) \cdot |\det(\vec{A})| &=
            \frac{|\det(\vec{A})|}{\sqrt{(2\pi)^n \cdot \det(\vec{A} \cdot \vec{A}^{\text{\,T}})}} \cdot 
            e^{-\frac{1}{2}(\vec{A}\cdot \vec{z} + \vec{\mu}-\vec{\mu})^{\text{T}} \cdot (\vec{A} \cdot \vec{A}^{\text{\,T}})^{-1} \cdot (\vec{A}\cdot \vec{z} +\vec{\mu}-\vec{\mu})} \\
            &= \frac{1}{\sqrt{(2\pi)^n}} \cdot e^{-\frac{1}{2}\, \vec{z}^{\text{\,T}} \cdot \vec{z}}
            \implies \vec{Z} \sim \mathcal{N}(0,\,\vec{I})
        .\end{align*}
        Abbiamo mostrato che $\vec{X} = \vec{A} \cdot \vec{Z} + \vec{\mu}$ è un vettore gaussiano di media $\vec{\mu}$ e matrice di covarianza $\vec{C}$
        \end{obsv}
        \begin{prty}
            Sia $\vec{X} = \vec{A} \cdot \vec{Z} + \vec{\mu}$ un vettore gaussiano $n$\nbdash dimensionale, 
            sia $\vec{C} = \vec{A} \cdot \vec{A}^{\text{\,T}}$ la matrice di covarianza di $\vec{X}$; 
            allora possiamo affermare che:
            \begin{enumerate}
                \item $c_{i,\,i} \in \vec{C}_{\vec{X}} \,:\, \begin{cases}
                    X_i \sim \mathcal{N}(\mu_i,\,c_{i,\,i}) & \text{se $c_{i,\,i} > 0$;} \\
                    P(X_i = \mu_i) = 1 & \text{se $c_{i,\,i} = 0$;}
                    \end{cases}$
                \item $\vec{G} \in \mathbb{R}^k \times \mathbb{R}^n \land \vec{h} \in \mathbb{R}^k
                    \implies \vec{Y} \coloneqq \vec{G} \cdot \vec{X} + \vec{h} \sim 
                    \mathcal{N}\left(\vec{G}\, \vec{\mu} + \vec{h},\, \vec{G}\, \vec{C}\, \vec{G}^{\text{\,T}}\right);$
                \item se $X_1,\, \ldots,\, X_{n}$ non sono correlate allora sono indipendenti.
            \end{enumerate}
        \end{prty}
        \begin{proof}
            \hfill
            \begin{enumerate}
                \item Dato che ogni componente $X_i$ del vettore aleatorio $\vec{X}$ si può esprimere come 
                    combinazione lineare di variabili aleatorie gaussiane indipendenti più una 
                    costante (dalla Definizione~\ref{defn:Vettore_gaussiano}); 
                    dall'Osservazione~\ref{obsv:Somma_gaussiane_indipendenti} possiamo dedurre che anche $X_i$ 
                    definita in tal modo è ancore una variabile aleatoria gaussiana.
                \item Scriviamo $\vec{Y}$ nel modo seguente: \[
                        \vec{G} \cdot \vec{X} + \vec{h} = \vec{G}(\vec{A} \cdot \vec{Z} + \vec{\mu}) + \vec{h}
                        = (\vec{G} \cdot \vec{A})\cdot \vec{Z} + (\vec{G} \cdot \vec{\mu} + \vec{h})
                    .\] Usando la Definizione~\ref{defn:Vettore_gaussiano} e 
                    l'Osservazione~\ref{obsv:Vettore_gaussiano} possiamo affermare che:
                    \begin{align*}
                        \text{E}(\vec{X}) = \vec{G}\, \vec{\mu} + \vec{h}, & &
                        \vec{C}_{\vec{X}} = \vec{G}\, \vec{C}\, \vec{G}^{\text{\,T}}
                    .\end{align*}
                \item Consideriamo una matrice di covarianza invertibile ($\vec{X}$ ha densità su 
                    $\mathbb{R}^n$); se le componenti $X_1,\, \ldots,\, X_{n}$ non sono correlate, allora 
                    $\vec{C}_{\vec{X}}$ è una matrice diagonale costruita come segue: \[
                        \vec{C}_{\vec{X}} = \text{diag}\left(\sigma_1^2 = \text{Var}(X_1),\, 
                        \ldots,\, \sigma_n^2 = \text{Var}(X_n)\right)
                    ;\] allora la densità di $\vec{X}$ si ottiene come segue:
                    \begin{align*}
                        f_{\vec{X}}(\vec{x}) &= \frac{1}{\sqrt{(2\pi)^n \cdot \det(\vec{C})}} \cdot e^{-\frac{1}{2}(\vec{x}-\vec{\mu})^{\text{T}}cd \vec{C}^{\,-1}\cdot(\vec{x}-\vec{\mu})} \\
                        &= \frac{1}{\sqrt{(2\pi)\cdot [\sigma_1^2\cdot \ldots\cdot \sigma_n^2]}} \cdot e^{-\frac{1}{2} \sum_{i=1}^{n} \left[\frac{x_i-\mu_i}{\sigma_i}\right]^2} \\
                        &= \prod_{i=1}^{n} \frac{1}{\sqrt{2\pi \sigma_i^2}} \cdot e^{-\frac{1}{2}\left[\frac{x_i-\mu_i}{\sigma_i}\right]^2} \\
                        &= \prod_{i=1}^{n} f_{X_i}(x_i)
                    .\end{align*}
                    Segue che le $X_i$ sono indipendenti. \qedhere
            \end{enumerate}
        \end{proof}
    \section{Distribuzione Gamma}
        \begin{defn}
            Una variabile aleatoria continua si dice avere distribuzione \emph{gamma} se la sua densità si ottiene dalla Definizione~\ref{defn:Densità_Gamma}:
            \begin{equation}\label{eq:Densità_Gamma_Eulero}
                f_{X}(x) = \begin{cases}
                    \frac{\lambda^\alpha}{\Gamma(\alpha)}\cdot x^{\alpha-1}\cdot e^{-\lambda\, x} & \text{se $x > 0$;} \\
                    0 & \text{se $x \leq 0$.}
                \end{cases}
            \end{equation}
            Nella precedente siano definiti i parametri $\alpha > 0,\, \lambda > 0$ e la funzione $\Gamma(\alpha)$ si chiama \emph{Gamma di Eulero} ed è definita come:
            \begin{align}\label{eq:Gamma_di_Eulero}
                \Gamma(\alpha) &\coloneqq \int_{0}^{\infty} \lambda^\alpha \cdot x^{\alpha -1}\cdot e^{-\lambda\, x}\, dx \notag\\
                &= \int_{0}^{\infty} y^{\alpha -1}\cdot e^{-y}\, dy & & [y=\lambda\cdot x]
            .\end{align}
        \end{defn}
        \begin{obsv}[Gamma di Eulero]
            Applicando la formula dell'integrazione per parti alla \eqref{eq:Gamma_di_Eulero}, nel caso in cui $\alpha > 1$, possiamo riscriverla come:
            \begin{align*}
                \Gamma(\alpha) &= \int_{0}^{\infty} y^{\alpha-1}\cdot e^{-y}\, dy 
                \overset{\scriptscriptstyle \text{P.P.} \textstyle}{=}
                \left[-y^{\alpha-1} \cdot e^{-y}\right]_{y=0}^{y=+\infty} + 
                \int_{0}^{\infty} (\alpha-1) \cdot y^{\alpha-2} \cdot e^{-y}\, dy \\
                &= (\alpha-1) \cdot \int_{0}^{\infty} y^{\alpha-2} \cdot e^{-y}\, dy \\
                &= (\alpha-1)\cdot \Gamma(\alpha-1)
            .\end{align*}
            Possiamo usare questo risultato per calcolare il valore della funzione Gamma sugli interi:
            \begin{align*}
                \Gamma(n) &= (n-1)\cdot \Gamma(n-1) \\
                &= (n-1)\cdot (n-2)\cdot \Gamma(n-2) \\
                &\;\; \vdots \\
                &= (n-1)! \cdot \Gamma(1)
            .\end{align*}
            Sapendo che per $n=1$ otteniamo: \[
                \Gamma(1) = \int_{0}^{\infty} e^{-y}\, dy = 1
            ,\] possiamo ottenere, per $n > 1,\, n \in \mathbb{N}$:
            \begin{equation}\label{eq:Gamma_di_Eulero_interi}
                \Gamma(n) = (n-1)!
            .\end{equation}
        \end{obsv}
        \begin{prty}[Momenti gamma]
            La generatrice dei momenti di una variabile aleatoria di tipo Gamma con parametri $\alpha,\, \lambda$ si ottiene come segue:
            \begin{align}\label{eq:Generatrice_momenti_Gamma}
                m_X(t) &\coloneqq \text{E}(e^{t\, X}) =
                \int_{0}^{\infty} e^{t\, X}\cdot \frac{\lambda^{\alpha}}{\Gamma(\alpha)}\cdot x^{\alpha-1}\cdot e^{-\lambda x}\, dx \notag\\
                &= \frac{\lambda^\alpha}{\Gamma(\alpha)}\cdot \int_{0}^{\infty} x^{\alpha-1} \cdot e^{-(\lambda-t)x}\, dx \notag\\
                &= \left(\frac{\lambda}{\lambda-t}\right)^\alpha \cdot 
                \frac{1}{\Gamma(\alpha)}\cdot 
                \underset{\Gamma(\alpha)}{\underbrace{\int_{0}^{\infty} y^{\alpha-1}\cdot e^{-y}\, dy}} 
                & & \left[y = (\lambda-t)x\right] \notag\\
                &= \left(\frac{\lambda}{\lambda-t}\right)^\alpha
            .\end{align}
        \end{prty}
        \begin{defn}
            Usando la funzione \eqref{eq:Generatrice_momenti_Gamma} possiamo ottenere valore atteso e varianza per una variabile aleatoria con distribuzione Gamma:
            \begin{itemize}
                \item $\text{E}(X) = m_X(0)^{\prime} = 
                    \frac{\alpha\cdot \lambda^{\alpha}}{(\lambda-0)^{\alpha+1}} = 
                    \frac{\alpha}{\lambda}$;
                \item $\text{E}(X^2) = m_X(0)^{\prime\prime} = 
                    \frac{\alpha(\alpha+1)\cdot \lambda^{\alpha}}{(\lambda-0)^{\alpha+2}} = 
                    \frac{\alpha(\alpha+1)}{\lambda^2}$;
                \item $\text{Var}(X) = \text{E}(X^2) - \text{E}(X)^2 = \frac{\alpha}{\lambda^2}$.
            \end{itemize}
        \end{defn}
        \begin{defn}[Somma di Gamma indipendenti]
            Consideriamo una successioni di variabili aleatorie indipendenti identicamente distribuite $X_1,\, \ldots,\, X_{n}$ tali che ciascuna di esse abbia densità Gamma  $X_i \sim \Gamma(\alpha_i,\, \lambda)$; allora la loro somma ha densità: \[
                X_1 + \ldots + X_{n} = \sum_{i=1}^{n} \left[X_i \sim \Gamma(\alpha_i,\, \lambda)\right] \sim \Gamma\left(\sum_{i=1}^n \alpha_i,\, \lambda\right)
            .\] Osservando che: \[
            \Gamma(1,\, \lambda) \sim \mathcal{E}(\lambda)
        ,\] possiamo affermare che, presa la successione $Y_1,\, \ldots,\, Y_{n}$ dove $Y_i \sim \mathcal{E}(\lambda)$, la distribuzione della somma delle $Y_i$ si ottiene come: \[
        Y_1 + \ldots + Y_{n} = \sum_{i=1}^n \left[Y_i \sim \mathcal{E}(1,\, \lambda)\right] \sim \Gamma(n,\, \lambda)
        .\] 
        \end{defn}
    \section{Distribuzioni derivate dalla Normale}
        \subsection{Distribuzione Chi-Quadro}
            \begin{defn}\label{defn:Distribuzione_chi_quadro}
                Prendiamo una successione $Z_1,\, \ldots,\, Z_{n}$ di variabili aleatorie gaussiane standard indipendenti; allora la somma dei loro quadrati $X \coloneqq Z^2_1 + \ldots + Z^2_{n}$ è una variabile aleatoria chiamata \emph{chi-quadro} con $n$ gradi di libertà, e si indica come: \[
                    X \sim \chi^2(n) \equiv \chi^2_n
                .\]
            \end{defn}
            \begin{obsv}[Riproducibilità]
                Se $X_1,\,X_2$ sono due variabili aleatorie con densità chi-quadro \underline{indipendenti}, rispettivamente a $n_1,\,n_2$ gradi di libertà, la loro somma ha la seguente densità: \[
                    W = X_1 + X_2 \sim \chi^2_{n_1} + \chi^2_{n_2} = \chi^2_{n_1+n_2}
                .\] Possiamo provare questo risultato usando direttamente la Definizione~\ref{defn:Distribuzione_chi_quadro}: sapendo che la somma di $X_1$ e $X_2$ è la somma di $n_1$ con $n_2$ normali standard indipendenti, è evidente che $W$ corrisponde a una chi-quadro con la somma dei gradi di libertà $n_1+n_2$.
            \end{obsv}
            \begin{prty}[Momenti chi-quadro]\label{prty:Generatrice_momenti_chi_quadro}
                Consideriamo il caso con $n=1$ gradi di libertà, e chiediamoci quanto vale la generatrice dei momenti di $X \sim \chi^2_1$, assumendo: \[
                    X = Z^2 \land Z \sim \mathcal{N}(0,\,1)
                ;\] allora otteniamo la generatrice dei momenti come:
                \begin{align*}
                    m_{X}(t) &= m_{Z^2}(t) =
                    \int_{-\infty}^{\infty} e^{t\, x^2}\cdot f_{Z}(x)\, dx = 
                    \int_{-\infty}^{\infty} e^{t\, x^2}\cdot \frac{1}{\sqrt{2\pi}}\, e^{-x^2 /2}\, dx \\
                    &= \int_{-\infty}^{\infty} \frac{1}{\sqrt{2\pi}}\cdot e^{-(1-2t)\frac{x^2}{2}}\, dx \\
                    &= (1-2t)^{-1 /2} \cdot \int_{-\infty}^{\infty} \underset{\sim\mathcal{N}(0,\,\bar{\sigma}^2)}{\underbrace{\frac{1}{\bar{\sigma}\, \sqrt{2\pi}} \cdot e^{-\frac{x^2}{2\bar{\sigma}^2}}}}\, dx & & \text{per } \bar{\sigma} = (1-2t)^{-1 /2} \\
                    &= (1-2t)^{-1 /2}
                .\end{align*}
                Nel caso più generale con $n$ gradi di libertà ($X \sim \chi^2(n)$), si ha:
                \begin{align}\label{eq:Generatrice_momenti_chi_quadro}
                    m_{X}(t) &= \text{E}(e^{t\,X}) = \text{E}(e^{t\,Z_1^2 + \ldots + t\,Z_n^2}) \notag\\
                             &= \text{E}(e^{t\,Z_1^2}\cdot \ldots \cdot e^{t\,Z_n^2}) \notag\\
                             &=\prod_{i=1}^{n} \text{E}(e^{t\,Z_i^2}) = (1-2t)^{-n /2} & & \text{per $Z_i$ indipendenti\notag}\\
                             &= \left(\frac{1}{1-2t}\right)^{n /2} = \left(\frac{1 /2}{1 /2 -t}\right)^{n /2} & & \text{dal caso $n=1$}
                .\end{align}
            \end{prty}
            \begin{defn}[Relazione tra gamma e chi-quadro]\label{defn:Relazione_gamma_chi_quadro}
                La densità Gamma di parametri $\alpha = n /2$ e $\lambda = 1 /2$  con $n \geq 1$ intero, è anche indicata come densità \emph{chi-quadro} con $n$ gradi di libertà.

                Possiamo provarlo direttamente dal risultato della Proprietà~\ref{prty:Generatrice_momenti_chi_quadro}: infatti la generatrice dei momenti della distribuzione chi-quadro ad $n$ gradi di libertà \eqref{eq:Generatrice_momenti_chi_quadro} corrisponde alla generatrice della distribuzione Gamma nel caso indicato sopra \eqref{eq:Generatrice_momenti_Gamma}.
            \end{defn}
        \subsection{Distribuzione \emph{t}}
            \begin{defn}
                Consideriamo le variabili aleatorie \underline{indipendenti} $Z \sim \mathcal{N}(0,\,1)$ e $C_n \sim \chi^2(n)$; allora la variabile aleatoria definita come segue:
                \begin{equation}\label{eq:Distribuzione_t}
                    T_n \coloneqq \frac{Z}{\sqrt{C_n /n}}
                .\end{equation}
                ha distribuzione chiamata \emph{t di Student} con $n$ gradi di libertà, e si indica con: \[
                    T_n \sim t(n)
            .\]
            \end{defn}
            \begin{prty}\label{prty:Distribuzione_t}
                Sia data una variabile aleatoria $T_n \sim t(n)$, allora possiamo affermare che:
                \begin{enumerate}
                    \item la distribuzione \emph{t} è simmetrica rispetto all'ascissa ($x=0$), ovvero $T_n \sim -T_n$;
                    \item per $n$ adeguatamente grande ($n \geq 100$) possiamo approssimare $T_n \sim \mathcal{N}(0,\,1)$;
                    \item per $n>2$ otteniamo $\text{E}(T_n)=0,\, \text{Var}(T_n)=\frac{n}{n-2}$.
                \end{enumerate}
            \end{prty}
            \begin{proof}
                \hfill
                \begin{enumerate}
                    \item Usando la definizione della distribuzione \emph{t} sappiamo che: \[
                        Z \sim \mathcal{N}(0,\,1) \sim -Z \implies
                        -T_n = \frac{-Z}{\sqrt{C_n /n}} \sim \frac{Z}{\sqrt{C_n /n}} = T_n
                    .\] 
                    \item Dalla funzione di ripartizione di $T_n$ otteniamo: \[
                            F_{T_n}(x) \coloneqq P(T_n \leq x) = P\left(\frac{Z}{\sqrt{C_n /n}} \leq x\right)
                        ;\] applicando il Teorema~\ref{thm:Legge_grandi_numeri_forte} possiamo affermare che, per $n$ grande, $X /n \rightarrow \text{E}(Z_i^2)=1$; quindi: \[
                        \lim_{n \to \infty} F_{T_n}(x) = 
                        \lim_{n \to \infty} P\left(\frac{Z}{\sqrt{C_n /n}} \leq x\right) \rightarrow
                        P\left(\frac{Z}{1} \leq x\right) = \Phi(x) = \frac{1}{\sqrt{2\pi}} \cdot \int_{-\infty}^{x} e^{-u^2 /2}\, du
                        .\qedhere\]
                \end{enumerate}
            \end{proof}
        \subsection{Distribuzione \emph{F}}
        \begin{defn}
            Consideriamo due variabili aleatorie indipendenti $C_n \sim \chi^2(n)$ e $C_m \sim \chi^2(m)$; allora la variabile aleatoria definita in modo seguente:
            \begin{equation}\label{eq:Distribuzione_F}
                F_{n,\,m} \coloneqq \frac{C_n /n}{C_m /m}
            ,\end{equation}
            ha una distribuzione chiamata \emph{F di Fisher} con $n$ e $m$ gradi di libertà, e la indichiamo con: \[
                F_{n,\,m} = \mathcal{F}(n,\,m)
            .\]
        \end{defn}
        \begin{prty}
            Sia data una variabile aleatoria con densità di Fisher $F \sim \mathcal{F}(n,\,m)$; allora possiamo affermare che $F \geq 0$, essendo quoziente di due variabili aleatorie con distribuzione chi-quadro (a loro volta somma di quadrati, quindi positive per definizione);
        \end{prty}
    \section{Quantili}
        \begin{defn}
            Chiamiamo \emph{quantile} di ordine $\alpha$, e lo indichiamo con $q_\alpha$ per 
            $\alpha \in [0,\,1]$, il valore ottenuto dalla funzione di ripartizione $F_W(w)$ di una 
            variabile aleatoria continua $W$ come: \[
                F_W(q_\alpha) = P(W \geq q_\alpha) = \alpha
            .\] La distribuzione viene suddivisa da $q_\alpha$ in due intervalli proporzionali ad $\alpha$ e a $(1-\alpha)$.
        \end{defn}
        \subsection{Quantili della normale standard}
            \begin{defn}\label{defn:Quantili_normale_standard}
                Data la variabile aleatoria $Z \sim \mathcal{N}(0,\,1)$ chiamiamo quantile di coda sinistra di ordine $\alpha$ il valore definito come: \[
                    P(Z > z_\alpha) = 1 - \Phi(z_\alpha) = \alpha
                .\]
            \end{defn}
            \begin{obsv}
                Una variabile aleatoria con distribuzione normale standard avrà probabilità inferiore o uguale a $z_\alpha$ nel $[100\cdot (1-\alpha)]\%$ dei casi. Inoltre usando la definizione appena enunciata possiamo scrivere: \[
                    z_\alpha \coloneqq \Phi^{-1}(1-\alpha)
                .\]
            \end{obsv}
        \subsection{Quantili della chi-quadro}
            \begin{defn}\label{defn:Quantili_chi_quadro}
                Data la variabile aleatoria $C \sim \chi^2(n)$ chiamiamo quantile di coda destra 
                di ordine $\alpha$ il valore definito come: \[
                    P(C > \chi^2_{\alpha,\,n}) = \alpha
                .\]
            \end{defn}
        \subsection{Quantili della \emph{t}}
            \begin{defn}\label{defn:Quantili_t_Student}
                Data la variabile aleatoria $T \sim t(n)$ chiamiamo quantile di coda destra 
                di ordine $\alpha$ il valore definito come: \[
                    P(T > t_{\alpha,\,n}) = \alpha
                .\]
            \end{defn}
            \begin{obsv}
                Data la simmetria della densità \emph{t} rispetto all'asse $x=0$, possiamo scrivere:
                \begin{align*}
                    \alpha &= P(-T \geq t_{\alpha,\,n})
                    = P(T \leq -t_{\alpha,\,n}) \\
                    &= 1 - P(T > -t_{\alpha,\,n})
                .\end{align*}
                Segue che:
                \begin{equation}\label{eq:Simmatria_quantile_distribuzione_t}
                    P(T \geq -t_{\alpha,\,n}) = 1 - \alpha \implies -t_{\alpha,\,n} = t_{1-\alpha,\,n}
                .\end{equation}
            \end{obsv}
        \subsection{Quantili della \emph{F}}
            \begin{defn}
                Data la variabile aleatoria $F \sim \mathcal{F}(n,\,m)$ chiamiamo quantile di coda destra 
                di ordine $\alpha$ il valore definito come: \[
                    P(F > f_{\alpha,\,n,\,m}) = \alpha
                .\]
            \end{defn}
            \begin{obsv}
                Analizziamo il caso di un quantile di ordine simmetrico ($1-\alpha$):
                \begin{align*}
                    \alpha &= P\left(\frac{C_n /n}{C_m /m} > f_{\alpha,\,n,\,m}\right) =
                    P\left(\frac{C_m /m}{C_n /n} < \frac{1}{f_{\alpha,\,n,\,m}}\right) \\
                           &= 1 - P\left(\frac{C_m /m}{C_n /n} \geq \frac{1}{f_{\alpha,\,n,\,m}}\right)
                .\end{align*}
                Segue che:
                \begin{equation}\label{eq:Simmetria_quantile_distribuzione_F}
                    P\left(\frac{C_m /m}{C_n /n} \geq f_{1-\alpha,\,n,\,m}\right) = 1-\alpha \implies
                    \frac{1}{f_{\alpha,\,n,\,m}} = f_{1-\alpha,\,n,\,m}
                .\end{equation}
            \end{obsv}


\end{document}
